\documentclass[10pt]{amsart}
\usepackage{amssymb,amsmath,amsthm,amsfonts}
\usepackage{graphicx,color,wrapfig,ytableau,mathtools}
\usepackage[foot]{amsaddr}
\usepackage{stmaryrd}
\usepackage{tikz,tikz-cd}
\usetikzlibrary{arrows.meta}
\usepackage{caption,enumerate}
\usepackage{todonotes}
\usepackage{hyperref,cleveref}
\usepackage[style=alphabetic,backend=biber]{biblatex}
\addbibresource{references.bib}

\newcommand\cred[1]{{\color{red} #1}}
\newcommand\cbl[1]{{\color{blue} #1}}

\newcommand{\Eugene}[1]{\todo[inline,color=red!30,author=Eugene]{#1}}
\newcommand{\Semeon}[1]{\todo[inline,color=blue!25,author=Semeon]{#1}}
\newcommand{\Daniil}[1]{\todo[inline,color=orange!60,author=Daniil]{#1}}

\newcommand{\sh}[2]{\langle #1,#2\rangle}

\allowdisplaybreaks

\def\tr{{\rm tr}\,}
\def\Tr{{\rm Tr}\,}
\newcommand{\Tnil}{\mathbb{T}^{\mathrm{nil}}}

\newtheorem{theorem}{Theorem}
\newtheorem{definition}[theorem]{Definition}
\newtheorem{lemma}[theorem]{Lemma}
\newtheorem{proposition}[theorem]{Proposition}
\newtheorem{assumption}[theorem]{Assumption}
\newtheorem{remark}[theorem]{Remark}
\newtheorem{corollary}[theorem]{Corollary}
\newtheorem{conjecture}[theorem]{Conjecture}
\newtheorem{example}[theorem]{Example}

\newtheorem{theoremA}{Theorem}
\renewcommand{\thetheoremA}{A}
\newtheorem{theoremB}{Theorem}
\renewcommand{\thetheoremB}{B}
\newtheorem{theoremC}{Theorem}
\renewcommand{\thetheoremC}{C}

\newcommand{\fosp}{{\mathfrak osp}}
\newcommand{\fg}{{\mathfrak g}}
\newcommand{\fsl}{{\mathfrak sl}}
\newcommand{\fb}{{\mathfrak b}}
\newcommand{\fp}{{\mathfrak p}}
\newcommand{\fh}{{\mathfrak h}}
\newcommand{\fn}{{\mathfrak n}}
\newcommand{\Dalpha}{D(2|1,\alpha)}
\newcommand{\Ddeg}{D(2|1,\mathrm{deg})}
\newcommand{\Daff}{D^{(1)}(2|1,\mathrm{deg})}
\newcommand{\spn}{\operatorname{span}}
\newcommand{\LInd}{\mathbf{L}^{(\bullet)}\widetilde{\Ind}}
\newcommand{\Lind}{\mathbf{L}\widetilde{\Ind}}
\newcommand{\Rcoind}{\mathbf{R}\widetilde{\coInd}}
\newcommand{\RcoInd}{\mathbf{R}^{(\bullet)}\widetilde{\coInd}}
\DeclareMathOperator{\Res}{Res}
\DeclareMathOperator{\Ind}{Ind}
\DeclareMathOperator{\ch}{ch}
\DeclareMathOperator{\coInd}{coInd}
\DeclareMathOperator{\Hom}{Hom}
\DeclareMathOperator{\U}{U}
\DeclareMathOperator{\Ext}{Ext}

% ---- colours (as on the blackboard) ----------------------------
\definecolor{evenroot}{RGB}{236, 72,153} % pink   : even roots
\definecolor{oddroot}{RGB}{250,204, 21}  % yellow : odd simple roots
\definecolor{oddtext}{RGB}{133, 77, 14}  % dark gold : odd labels
\definecolor{eventext}{RGB}{190, 24,93}  % dark pink : even labels
\definecolor{negmark}{RGB}{ 56,189,248}  % cyan : circled-minus marks
\definecolor{oddblue}{RGB}{ 59,130,246}  % blue : non-simple odd roots
\definecolor{oddbluetext}{RGB}{30, 64,175} % dark blue : their labels

\tikzset{
	rootfig/.style={
		x={(1.2cm,0cm)}, y={(0.54cm,0.66cm)}, z={(0cm,1.2cm)},
		line cap=round, line join=round,
		lab/.style={font=\small},
		wlab/.style={lab, anchor=west,        xshift= 5pt},
		elab/.style={lab, anchor=east,        xshift=-5pt},
		nwlab/.style={lab, anchor=north west, xshift= 4pt, yshift=-3pt},
		nelab/.style={lab, anchor=north east, xshift=-4pt, yshift=-3pt},
		selab/.style={lab, anchor=south east, xshift=-4pt, yshift= 3pt},
		slab/.style={lab, anchor=south,       yshift= 3pt},
		nlab/.style={lab, anchor=north,       yshift=-3pt},
		axis/.style={evenroot!70, line width=0.7pt},
		axistip/.style={axis, -{Latex[length=2.8mm,width=2.2mm]}},
		axishidden/.style={axis, dash pattern=on 2pt off 2.4pt},
		hidden/.style={black!45, dash pattern=on 2.5pt off 2pt},
		solid edge/.style={black!75, line width=0.6pt},
	}
}

%%%%%%%%%%%%%%%%%%%%%%%%%%%%%%%%%%%%%%%%%%%%%%%%%%%%%%%%%%%%%%%%%%%%%%%%
%%%%%%%%%               SPACE FILLING SETTINGS               %%%%%%%%%%%
%%%%%%%%%%%%%%%%%%%%%%%%%%%%%%%%%%%%%%%%%%%%%%%%%%%%%%%%%%%%%%%%%%%%%%%%
\topmargin -1in
\headheight 0in
\headsep 0in
\textheight 9.5in
\textwidth 7.1in
\oddsidemargin 0in
\evensidemargin 0in
\headheight 24pt
\headsep 0.25in
\hoffset=-0.8cm
\voffset=0.2cm
%%%%%%%%%%%%%%%%%%%%%%%%%%%%%%%%%%%%%%%%%%%%%%%%%%%%%%%%%%%%%%%%%%%%%%%%

% --- Affiliations ---
\newcommand{\affil}[1]{\textsuperscript{#1}}
\makeatletter
\newcommand{\affiltext}[2]{
  \begingroup
  \renewcommand{\thefootnote}{#1}
  \phantomsection
  \footnotetext{#2}
  \endgroup
}
\makeatother

\title{Categorification of the genus two DAHA}
\author{Semeon~Arthamonov\affil{1}}
\author{Ievgen~Makedonskyi\affil{1}}
\author{Daniil~Sarafannikov\affil{2}}

\begin{document}
	
	\begin{abstract}
		We construct three derived endofunctors on the derived category of graded integrable representations of the current algebra of a flat degeneration of $D(2|1,\alpha)$, and prove that their classes in the Grothendieck group are the genus two Macdonald operators. Computing the graded $\Ext$ pairing explicitly, we deduce the self-adjointness of the genus two Macdonald operators and the orthogonality of genus $2$ Macdonald polynomials with respect to the certain skew-bilinear form.
	\end{abstract}
	
	\maketitle
	
    \affiltext{1}{Beijing Institute of Mathematical Sciences and Applications (BIMSA), Beijing, China. \textit{Email:} arthamonov@bimsa.cn, mak@bimsa.cn}
    \affiltext{2}{Qiuzhen College, Tsinghua University, Beijing, China. \textit{Email:} sardanis123@gmail.com}
	
	\section{Introduction}
	
	The double affine Hecke algebra of Cherednik \cite{Cherednik-2005} and its spherical subalgebra provide the algebraic framework for the theory of Macdonald polynomials. In the rank one case the spherical DAHA of type $A_1$ acts on the ring of symmetric Laurent polynomials in one variable, the Macdonald difference operator is one of its generators and the Macdonald polynomials are the eigenfunctions of this operator. From the topological point of view this algebra is attached to the once-punctured torus: its mapping class group $SL(2,\mathbb{Z})$ acts on the algebra by automorphisms.
	
	In \cite{ArthamonovShakirov-2019} a genus two counterpart of this picture was proposed. The construction of \cite{ArthamonovShakirov-2019} starts from three commuting $q$-difference operators $\hat O_{A_{12}}, \hat O_{A_{13}}, \hat O_{A_{23}}$ in three variables $X_1,X_2,X_3$ depending on two parameters $q$ and $t$. The algebra generated by these operators together with the operators of multiplication by $X_k+X_k^{-1}$ carries an action of the Dehn twists of the closed genus two surface by automorphisms, and these automorphisms satisfy all the relations of the mapping class group of that surface. The joint eigenfunctions of the three commuting operators are the genus two analogue of the $A_1$ Macdonald polynomials. Following \cite{ArthamonovShakirov-2019} we call $\hat O_{A_{ij}}$ the {\it genus two Macdonald operators}, the label $A_{ij}$ refers to a curve on the genus two surface.
	
	The purpose of the present paper is to categorify these operators, that is, to realize them as the classes of a family of derived functors between categories of representations. The mechanism we use is similar to one which was used in \cite{KKhM}: one takes a Lie superalgebra with a suitable family of parabolic subalgebras, considers the induction, restriction and coinduction functors between the corresponding categories of graded representations, and computes their action on the Grothendieck groups, which are identified with rings of Laurent polynomials by means of characters. The similar categorification was made in \cite{bezrukavnikov2012affine} using the geometric approach. It is interesting question to get such a geometric analogue of our results.
	
	The algebra in question is a degeneration of the affinization of the exceptional family of simple Lie superalgebras $D(2|1,\alpha)$. Recall that $D(2|1,\alpha)$ has dimension $(9|8)$, its even part is $\fsl_2\oplus\fsl_2\oplus\fsl_2$ and its rank is three. The whole dependence on the parameter $\alpha$ is concentrated in the Lie bracket component $S^2\fg_{\overline 1}\rightarrow \fg_{\overline 0}$ (as map of vector spaces), so that setting this component of the Lie bracket to zero produces a Lie superalgebra $\Ddeg$ which does not depend on $\alpha$. It is isomorphic to $D(2|1,0,0,0)$, see Definition~\ref{def:degeneratedLieSuperalgebra}. We denote by $\Daff$ the affinization of $\Ddeg$. It has four simple roots $\alpha_0,\alpha_1,\alpha_2,\alpha_3$, all of them odd.
	
	The degeneration is what makes the algebra convenient to work with, and it is also what creates the additional grading we need. Namely, $\Daff$ carries three compatible gradings: the weight grading with respect to the Cartan subalgebra, with lattice $\mathbb{Z}^3$, the grading by the degree in the loop variable $z$, and the $\mathbb{Z}$-grading in which the even part sits in degree $0$ and the odd part in degree $1$, which exists precisely because the bracket of two odd elements vanishes. We work with the representations graded with respect to all three gradings, and we denote the corresponding formal variables by $X_1,X_2,X_3$, by $q$ and by $-t$. The variables $X_k$ are naturally labelled by pairs: $X_k=X^{(\alpha_i+\alpha_j)/2}$ for $\{i,j,k\}=\{1,2,3\}$, which is exactly the labelling of the variables of \cite{ArthamonovShakirov-2019}. The two parameters $q$ and $t$ of \cite{ArthamonovShakirov-2019} thus acquire a representation theoretic meaning: $q$ is the loop grading and $-t$ is the odd grading.
	
	The parabolic subalgebras of a Kac--Moody superalgebra are defined by choosing a subset of the simple roots and adjoining the corresponding negative root vectors to the Borel subalgebra. In the degenerate case this recipe is not available, and we give a direct construction of the degenerate parabolic subalgebras $\fp_J$, $J\subset\{0,1,2,3\}$, in Definition~\ref{def:parabolicSubalgebras}. The main object of the paper is the parabolic subalgebra
	\[\fp_{123}\simeq \Ddeg\otimes\Bbbk[z]\]
	together with its three maximal parabolic subalgebras $\fp_{12},\fp_{13},\fp_{23}$. The restriction functor $\Res_{\fp_{ij}}^{\fp_{123}}$ has a left and a right adjoint, the integrable induction and coinduction functors:
	\begin{equation*}
		\begin{tikzcd}[column sep=4.2cm]
			\mathcal{O}^{gr}(\fp_{123})
			\arrow[r, "{\Res_{\fp_{ij}}^{\fp_{123}}}" description]
			& \mathcal{O}^{gr}(\fp_{ij})
			\arrow[l, bend right=17, "{\widetilde{\Ind}_{\fp_{ij}}^{\fp_{123}}}"']
			\arrow[l, bend left=17, "{\widetilde{\coInd}_{\fp_{ij}}^{\fp_{123}}}"]
		\end{tikzcd}
	\end{equation*}
	
	The first main property of this pair of algebras is that the induction and the coinduction differ only by a shift of the gradings.
	
	\begin{theoremA}
		\begin{equation}\label{eq:IndCoind}
			{\Rcoind}_{\fp_{ij}}^{\fp_{123}}\simeq {\Lind}_{\fp_{ij}}^{\fp_{123}}\sh{0}{-2}[-2].
		\end{equation}
	\end{theoremA}
	
	The same property holds for the pairs $\fb\otimes\Bbbk[\xi]\subset\fp_i\otimes\Bbbk[\xi]$, where $\fb$ and $\fp_i$ are a Borel and a minimal parabolic subalgebra of a semisimple Lie algebra and $\xi$ is an odd variable. It would be interesting to describe all the pairs of Lie superalgebras with this property, we do not know how large this class is.
	
	The Lie superalgebra $D^{(1)}(2|1,\alpha)$ carries an action of the symmetric group $\mathfrak{S}_4$ permuting the four simple roots. On the non-degenerate algebra these maps change the parameter $\alpha$ and hence are not automorphisms. After the degeneration the dependence on $\alpha$ disappears, and we obtain an action of $\mathfrak{S}_4$ on $\Daff$ by automorphisms. Twisting by an automorphism is an exact functor, and the automorphism $\pi_{(i,j)(k,0)}$ preserves the parabolic subalgebra $\fp_{ij}$, so it gives an exact endofunctor of $\mathcal{O}^{gr}(\fp_{ij})$. The functors we study are the compositions
	\[\mathcal{F}_{ij}:=\Lind_{\fp_{ij}}^{\fp_{123}} \circ \pi_{(i,j)(k,0)}\circ {\Res}_{\fp_{ij}}^{\fp_{123}},\]
	see Definition~\ref{def:Fdefinition}. Our main categorification result is the following.
	
	\begin{theoremB}
		\[[\mathcal{F}_{ij}]=t\, \hat O_{A_{ij}}.\]
	\end{theoremB}
	
	Thus the functors $\mathcal{F}_{ij}$ categorify the genus two Macdonald operators of \cite{ArthamonovShakirov-2019}. It is natural to ask whether the relations satisfied by these operators can be lifted to isomorphisms of functors. The commutativity of the operators $\hat O_{A_{ij}}$ is proved in \cite{ArthamonovShakirov-2019}, so that the classes $[\mathcal{F}_{ij}]$ commute. We expect this to hold on the categorical level.
	
	\begin{conjecture}
		For any $i,j,k$ there is an isomorphism of functors
		\[\mathcal{F}_{ij}\circ\mathcal{F}_{ik}\simeq \mathcal{F}_{ik}\circ \mathcal{F}_{ij}.\]
	\end{conjecture}
	
	Categorification gives more than the operators themselves. The category
	$\mathcal{O}^{gr}(\fp_{123})$ carries a natural graded $\Ext$ pairing
	\[(M,N)_{\Ext}:=\sum_{i=0}^{\infty}(-1)^i\dim_{q,t}\bigl(\Ext^{i}_{gr}(M ,N)\bigr),\]
	see Subsection~\ref{ssec:ExtPairing}, and the adjunctions between the functors above turn into a statement about the operators $\hat O_{A_{ij}}$. This form is additive with respect to short exact sequences in each argument, hence descends to a pairing on the Grothendieck group, and it is skew-linear in the second argument with respect to the substitution $f^\star=f|_{t\mapsto t^{-1},q\mapsto q^{-1}}$. Identifying $K_0(\mathcal{O}^{gr}(\fp_{123}))$ with a ring of symmetric Laurent polynomials by means of characters, this pairing takes the explicit form
	\[(f,g)=\bigl(f\,\Delta\, g^\star\bigr)_0,\qquad \Delta:=(q;q)^3\frac{\prod_{l \geq 0}\prod_{i=1}^3(1-X_i^2 q^l)(1-X_i^{-2} q^l)}{\prod_{l \geq 0}\prod_{a,b,c\in\{\pm 1\}}(1-X_1^aX_2^bX_3^c q^lt)},\]
	where $(\cdot)_0$ denotes the constant term with respect to the variables $X_i$, see Lemma~\ref{lem:ExtPairingExplicit}. The density $\Delta$ is the genus two counterpart of the Macdonald weight function: the numerator is the contribution of the even part of the algebra and the denominator is the contribution of the odd part.
	
	The property \eqref{eq:IndCoind} says exactly that a suitably normalized version of $\mathcal{F}_{ij}$ is self-adjoint with respect to this pairing. Set
	\[\mathcal{G}_{ij}:=\mathcal{F}_{ij}\sh{0}{-1}[-1].\]
	Then
	\[\Hom_{gr}(\mathcal{G}_{ij}M,N)\simeq \Hom_{gr}(M,\mathcal{G}_{ij}N),\qquad [\mathcal{G}_{ij}]=\hat O_{A_{ij}},\]
	and passing to the Grothendieck group we obtain the following property of the genus two Macdonald operators.
	
	\begin{theoremC}
		\[\bigl(\hat O_{A_{ij}}f\cdot \Delta\, g^{\star}\bigr)_0
		=\bigl(f\cdot \Delta\,(\hat O_{A_{ij}}g)^{\star}\bigr)_0.\]
	\end{theoremC}
	
	
	As a consequence, the joint eigenfunctions of the three operators $\hat O_{A_{12}},\hat O_{A_{13}},\hat O_{A_{23}}$ are pairwise orthogonal with respect to $\Delta$. These eigenfunctions were constructed in \cite{ArthamonovShakirov-2019}, where they are shown to be indexed by triples $\mathbf{i}$ of non-negative integers and to specialize to the $A_1$ Macdonald polynomials, we denote them by $\Psi_{\mathbf{i}}$ and refer to \cite{ArthamonovShakirov-2019} for their construction and their basic properties. For $\mathbf{i}\neq\mathbf{j}$ one has
	\[\bigl( \Psi_{\mathbf{j}}\,\Delta\, \Psi^{\star}_{\mathbf{i}} \bigr)_0=0,\]
	see Corollary~\ref{cor:Orthogonality}. This is the genus two analogue of the orthogonality of Macdonald polynomials with respect to the Macdonald weight.
	
	The paper is organized as follows. In Section~\ref{sec:LieSuperalgebras} we recall the generalities on Lie superalgebras, their root systems and parabolic subalgebras, and then introduce $\Ddeg$, its affinization $\Daff$, the parabolic subalgebras $\fp_J$ and the diagrammatic automorphisms. Section~\ref{sec:Categories} contains the homological preliminaries, the definition of the categories of representations we work with, the computation of their Grothendieck groups in terms of characters and the definition of the $\Ext$ pairing. Section~\ref{sec:Functors} deals with the induction, restriction and coinduction functors, first in general and then for the parabolic subalgebras of $\Daff$, the main computation there is Lemma~\ref{lem:InductionOperator}. Section~\ref{sec:Categorification} recalls the genus two Macdonald operators of \cite{ArthamonovShakirov-2019} and contains the categorification theorem. Section~\ref{sec:Orthogonality} contains the comparison \eqref{eq:IndCoind} of induction and coinduction and the orthogonality statements. %Section~\ref{sec:Discussion} collects open questions.
	
	\section{Lie superalgebras}\label{sec:LieSuperalgebras}
	
	Throughout the paper $\Bbbk$ denotes an algebraically closed field of characteristic zero. Unless stated otherwise, all the vector spaces, algebras and tensor products are taken over $\Bbbk$.
	
	In this section we first recall the general notions concerning Lie superalgebras, their root systems and parabolic subalgebras. We then specialize to the Lie superalgebra which is the main object of this paper, namely to a flat degeneration of $\Dalpha$ and to its affinization. We do not prove here any of the general statements about Lie superalgebras, the standard references are \cite{Mus} and \cite{DictionaryOn}.
	
	\subsection{Generalities}
	Recall that a Lie superalgebra is a $\mathbb{Z}/2\mathbb{Z}$-graded $\Bbbk$-linear space $\fg = \fg_{\overline{0}}\oplus\fg_{\overline{1}}$ together with a bilinear map $\left[\ ,\ \right]:\fg\otimes\fg\rightarrow\fg$ such that $\left[\fg_{\overline{i}}, \fg_{\overline{j}}\right]\subset\fg_{\overline{i+j}}$ and such that for all homogeneous elements $X, Y, Z \in \fg$ the following identities hold:
	\begin{align*}
		\left[X,Y\right]=-(-1)^{\deg{X}\deg{Y}}\left[Y,X\right],\\
		\left[X,\left[Y,Z\right]\right]=\left[\left[X,Y\right],Z\right]+(-1)^{\deg{X}\deg{Y}}\left[Y,\left[X,Z\right]\right].
	\end{align*}
	Here and below $\deg X \in \mathbb{Z}/2\mathbb{Z}$ denotes the parity of a homogeneous element $X$, and all the identities involving parities are stated for homogeneous elements and extended by linearity.
	
	In particular $\fg_{\overline 0}$ is an ordinary Lie algebra and $\fg_{\overline 1}$ is a $\fg_{\overline 0}$-module. A module over a Lie superalgebra is always assumed to be $\mathbb{Z}/2\mathbb{Z}$-graded and the action is assumed to respect this grading. We denote by $\U(\fg)$ the universal enveloping algebra of $\fg$. It is an associative superalgebra and the Poincar\'e--Birkhoff--Witt theorem holds for it in the form
	\[\U(\fg)\simeq \U(\fg_{\overline 0})\otimes \Lambda(\fg_{\overline 1})\]
	as $\mathbb{Z}/2\mathbb{Z}$-graded vector spaces, see \cite[Chapter~6]{Mus}.
	
	Throughout the paper we deal with Lie superalgebras which admit a semidirect decomposition
	\begin{equation}\label{eq:ReductivePart}
		\fg \simeq \fg^{red}\ltimes R,
	\end{equation}
	where $R$ is a (pro-)nilpotent ideal and $\fg^{red}$ is a reductive Lie algebra contained in $\fg_{\overline 0}$. We call $\fg^{red}$ the {\it reductive part} of $\fg$ and $R$ the {\it radical} of $\fg$. This decomposition is a part of the data and is fixed once and for all for each algebra we consider.
	
	\subsection{Root systems, Borel and parabolic subalgebras}
	Let $\fh \subset \fg_{\overline 0}$ be a Cartan subalgebra of the reductive part $\fg^{red}$ and assume that $\fh$ acts semisimply on $\fg$ by the adjoint action. Then we have the root decomposition
	\[\fg=\fh \oplus \bigoplus_{\gamma \in \Phi}\fg_{\gamma},\qquad \fg_\gamma:=\{X \in \fg \mid [H,X]=\gamma(H)X ~\text{for all}~ H \in \fh\},\]
	and the root system $\Phi \subset \fh^*$ splits into the even and the odd part
	\[\Phi=\Phi_{\overline 0}\sqcup \Phi_{\overline 1},\qquad \Phi_{\overline i}:=\{\gamma \in \Phi\mid \fg_\gamma \cap \fg_{\overline i}\neq 0\}.\]
	A choice of a Borel subalgebra $\fb \supset \fh$ is the same thing as a choice of a subset $\Phi_+\subset \Phi$ of positive roots together with the corresponding subset of simple roots.
	
	In contrast with the case of semisimple Lie algebras, for a Lie superalgebra different systems of simple roots need not be conjugate to each other, so that a Lie superalgebra has in general several non-equivalent Dynkin diagrams, see \cite[Chapter~3]{DictionaryOn} and \cite{DrinfieldSecond}. In this paper we do not need this flexibility: we fix one particular system of simple roots, described in Subsection~\ref{ssec:Ddeg} below, and all the constructions of the paper refer to this fixed choice.
	
	\begin{definition}\label{def:WeylGroupConvention}
		Throughout the paper the {\it Weyl group} $W$ of a Lie superalgebra $\fg$ means the Weyl group of the root system $\Phi_{\overline 0}$ of the reductive part $\fg^{red}$, that is, the group generated by the reflections $s_\gamma$, $\gamma \in \Phi_{\overline 0}$.
	\end{definition}
	
	Note that $W$ acts on $\fh^*$ and preserves $\Phi$, but it does not act transitively on the set of the systems of simple roots. All the notions such as dominance of a weight and integrability of a module refer to this group.
	
	Finally, we use the following notational convention for parabolic subalgebras: for a subset $J$ of the set of indices of the simple roots we write $\fp_J$, always omitting the braces, so that $\fp_{ij}:=\fp_{\{i,j\}}$ and $\fp_{123}:=\fp_{\{1,2,3\}}$. The precise definition of $\fp_J$ in the situation we need is given in Definition~\ref{def:parabolicSubalgebras}. As explained in Remark~\ref{rem:ParabolicDegeneration}, the naive definition used in the semisimple case is not appropriate here.
	
	\subsection{The Lie superalgebra \texorpdfstring{$\Ddeg$}{Ddeg}}\label{ssec:Ddeg}
	Consider the Lie superalgebra $\fg=\fg_{\overline 0}\oplus \fg_{\overline 1}=\Dalpha$ of dimension $(9|8)$. Its even part is isomorphic to $\fsl_2 \oplus \fsl_2 \oplus \fsl_2$.
	The weight lattice of this algebra is isomorphic to $\mathbb{Z}^3$. We denote by $\varepsilon_1, \varepsilon_2,\varepsilon_3$ the natural basis of this lattice. The positive even roots of this algebra are $2 \varepsilon_i$, $i=1,2,3$.
	
	The odd part $\fg_{\overline 1}$ is isomorphic to the tensor product of three fundamental representations, $V_{\varepsilon_1}\otimes V_{\varepsilon_2}\otimes V_{\varepsilon_3}$.
	
	For generic values of the parameter $\alpha$ the Lie superalgebra $\Dalpha$ is simple. In this paper we deal with a flat degeneration of this algebra, which we denote by $\Ddeg$. To define it, introduce a parameter $\tau$ and consider the flat family of multiplications on the underlying superspace
	\begin{equation*}
		[X,Y]_\tau:=\begin{cases}
			[X,Y], &\text{ if }X \in  \Dalpha_{\overline 0}\text{ or }Y \in  \Dalpha_{\overline 0},\\
			\tau [X,Y], &\text{ if }X,Y \in  \Dalpha_{\overline 1}.
		\end{cases}
	\end{equation*}
	One checks directly that $[\cdot,\cdot]_\tau$ is a Lie superbracket, i.e.\ that it satisfies the Jacobi identity and super-antisymmetry.
	For $\tau \neq 0$ this bracket is isomorphic to the original one, so it defines a simple Lie superalgebra. The case of interest for us is the special point $\tau=0$.
	\begin{definition}\label{def:degeneratedLieSuperalgebra}
		$\Ddeg$ is the Lie superalgebra which is isomorphic to $\Dalpha$ as a graded vector space and is equipped with the super Lie bracket $[\cdot,\cdot]_0$.
	\end{definition}
	\begin{remark}
		Note that the degenerate Lie superalgebra $\Ddeg$ is in fact isomorphic to $D(2|1,0,0,0)$.
	\end{remark}
	
	This Lie superalgebra carries a natural $\mathbb{Z}$-grading in which $\fg_{\overline 0}$ sits in degree zero and $\fg_{\overline 1}$ sits in degree $1$. Note that the weights of the basis elements of $\fg_{\overline 1}$
	are $\pm \varepsilon_1 \pm \varepsilon_2 \pm \varepsilon_3$. Introduce the notation
	\[\alpha_1=-\varepsilon_1+\varepsilon_2 +\varepsilon_3,\]
	\[\alpha_2=\varepsilon_1-\varepsilon_2 +\varepsilon_3,\]
	\[\alpha_3=\varepsilon_1+\varepsilon_2 -\varepsilon_3.\]
	We choose these roots to be the simple positive roots. In particular all the simple roots of $\Ddeg$ are odd. This is the system of simple roots which is fixed once and for all throughout the paper. We also write $\theta:=\varepsilon_1+\varepsilon_2+\varepsilon_3=\alpha_1+\alpha_2+\alpha_3$ for the highest root. The whole root system is shown on Figure~\ref{fig:RootSystem}.
	\begin{figure}[ht]
		\centering
		\begin{tikzpicture}[rootfig]
			% ---------------- title -------------------------------------------
			% ------- axes: solid outside the cube, dashed inside -------------
			% the arrow tips are put on the outer ends only, so each half axis
			% is drawn from the face of the cube outwards
			\draw[axistip] (-1,0,0) -- (-2.6,0,0);
			\draw[axistip] ( 1,0,0) -- ( 2.6,0,0);
			\draw[axistip] (0,-1,0) -- (0,-3.5,0);
			\draw[axistip] (0, 1,0) -- (0, 3.5,0);
			\draw[axistip] (0,0,-1) -- (0,0,-2.6);
			\draw[axistip] (0,0, 1) -- (0,0, 2.6);
			% the parts running through the interior of the cube
			\draw[axishidden] (-1,0,0) -- (1,0,0)
			(0,-1,0) -- (0,1,0)
			(0,0,-1) -- (0,0,1);
			
			\fill[black!70] (0,0,0) circle (0.8pt);
			
			% ------- the cube ------------------------------------------------
			\draw[hidden]
			(-1,1,-1) -- ( 1,1,-1)
			(-1,1,-1) -- (-1,-1,-1)
			(-1,1,-1) -- (-1, 1, 1);
			\draw[solid edge]
			(-1,-1,-1) -- (1,-1,-1) -- (1,-1,1) -- (-1,-1,1) -- cycle
			(-1,-1, 1) -- (-1,1,1) -- (1,1,1) -- (1,-1,1)
			( 1,-1,-1) -- ( 1,1,-1) -- (1,1,1);
			
			% ------- even roots ----------------------------------------------
			\foreach \p in {(2,0,0),(-2,0,0),(0,2,0),(0,-2,0),(0,0,2),(0,0,-2)}
			\filldraw[fill=evenroot, draw=black!60, line width=0.5pt] \p circle (2.2pt);
			
			\node[eventext, slab]  at ( 2,0,0) {$2\varepsilon_1$};
			\node[eventext, nlab]  at (-2,0,0) {$-2\varepsilon_1$};
			\node[eventext, selab, xshift=2pt, yshift=2pt]   at (0, 2,0) {$2\varepsilon_2$};
			\node[eventext, nwlab, xshift=-8pt, yshift=-2pt] at (0,-2,0) {$-2\varepsilon_2$};
			\node[eventext, wlab]  at (0,0, 2) {$2\varepsilon_3$};
			\node[eventext, wlab]  at (0,0,-2) {$-2\varepsilon_3$};
			
			% ------- odd roots -----------------------------------------------
			\coordinate (ppp) at ( 1, 1, 1);
			\coordinate (ppm) at ( 1, 1,-1);
			\coordinate (pmp) at ( 1,-1, 1);
			\coordinate (pmm) at ( 1,-1,-1);
			\coordinate (mpp) at (-1, 1, 1);
			\coordinate (mpm) at (-1, 1,-1);
			\coordinate (mmp) at (-1,-1, 1);
			\coordinate (mmm) at (-1,-1,-1);
			
			\foreach \p in {mpp,pmp,ppm}
			\filldraw[fill=oddroot, draw=black!60, line width=0.4pt] (\p) circle (2.2pt);
			\foreach \p in {ppp,pmm,mpm,mmp,mmm}
			\filldraw[fill=oddblue, draw=black!60, line width=0.4pt] (\p) circle (2.2pt);
			
			\node[oddbluetext, wlab]  at (ppp) {$\theta$};
			\node[oddtext, wlab]      at (ppm) {$\alpha_3$};
			\node[oddtext, nwlab]     at (pmp) {$\alpha_2$};
			\node[oddbluetext, wlab]  at (pmm) {$-\alpha_1$};
			\node[oddtext, slab]      at (mpp) {$\alpha_1$};
			\node[oddbluetext, elab]  at (mpm) {$-\alpha_2$};
			\node[oddbluetext, elab]  at (mmp) {$-\alpha_3$};
			\node[oddbluetext, nelab, xshift=-6pt] at (mmm) {$-\theta$};
			
			% ------- legend ---------------------------------------------------
			\node[anchor=west, align=left, font=\small]
			at (canvas cs:x=3.6cm,y=1.1cm)
			{{\color{oddtext}\textbf{odd simple roots}}\\[1pt]
				{\color{oddtext}$\alpha_1,\alpha_2,\alpha_3$}\\[5pt]
				{\color{oddbluetext}\textbf{remaining odd roots}}\\[1pt]
				{\color{oddbluetext}$\pm\theta,\ -\alpha_1,\ -\alpha_2,\ -\alpha_3$}\\[5pt]
				{\color{eventext}\textbf{even roots}}\\[1pt]
				{\color{eventext}$\pm2\varepsilon_1,\ \pm2\varepsilon_2,\ \pm2\varepsilon_3$}\\[5pt]
				$\theta=\varepsilon_1+\varepsilon_2+\varepsilon_3$\\
				$\phantom{\theta}=\alpha_1+\alpha_2+\alpha_3$\\[2pt]
				$\alpha_0=\delta-\theta$};
		\end{tikzpicture}
		\caption{The root system of $\Ddeg$. The eight odd roots are the vertices of a cube and the six even roots $\pm2\varepsilon_i$ lie on the coordinate axes. The three odd simple roots are marked in yellow. The vertex $-\theta$ is the finite part of the simple affine root $\alpha_0=\delta-\theta$ of $\Daff$.}
		\label{fig:RootSystem}
	\end{figure}
	Then $2 \varepsilon_1 =\alpha_2+\alpha_3$, and similarly for the other two.
	
	We fix once and for all the following notation for the formal exponents of the weights. For $\gamma \in \fh^*$ we write $X^{\gamma}$ and set
	\[X_i:=X^{\varepsilon_i},\quad i=1,2,3,\qquad q:=X^{\delta},\]
	where $\delta$ is the imaginary root of the affine algebra introduced in Subsection~\ref{ssec:Affinization}. Thus
	\[X^{\alpha_1}=X_1^{-1}X_2X_3,\quad X^{\alpha_2}=X_1X_2^{-1}X_3,\quad X^{\alpha_3}=X_1X_2X_3^{-1},\]
	and for $\{i,j,k\}=\{1,2,3\}$ one has
	\begin{equation*}\label{eq:XviaAlpha}
		X_{ij}=X_{k}=X^{(\alpha_i+\alpha_j)/2}.
	\end{equation*}
	The variable $t$ plays a different role: it does not come from the weight lattice, but records the $\mathbb{Z}$-grading of $\Ddeg$ described above, i.e.\ the odd degree. In particular $X^{\gamma}$ never contains $t$, and the factors $t$ appearing in the characters below are written separately.
	
	For the Lie algebra $\fsl_2$ we use the standard basis $\{e,h,f\}$.
	Consider the basis of the irreducible representation $V_\varepsilon=\spn\{v_+,v_-\}$ such that $v_+$ is the highest weight vector, $v_-$ is the lowest weight vector, $e v_-=v_+$ and $f v_+ =v_-$. We add lower indices to the elements of our Lie superalgebra, i.e.\ the basis of $\Ddeg_{\overline 0}$ is $\{e_i, f_i, h_i\}$, $i=1,2,3$.
	We choose the even root vectors as
	\[e_{2 \varepsilon_i}:=e_{i},~ e_{-2 \varepsilon_i}:=f_{i},\quad i=1,2,3,\]
	and the odd root vectors as
	\[e_{\alpha_1}:=v_- \otimes v_+ \otimes v_+, ~e_{\alpha_2}:=v_+ \otimes v_- \otimes v_+, ~e_{\alpha_3}:=v_+ \otimes v_+ \otimes v_-.\]
	
	We denote by $\fb$ the Borel subalgebra spanned by the Cartan subalgebra and the positive root vectors, i.e.
	\[\fb =\fh \oplus \spn \langle e_{\alpha_1}, e_{\alpha_2}, e_{\alpha_3},\ 
	e_{2\varepsilon_1}, e_{2\varepsilon_2}, e_{2\varepsilon_3},\ e_{\theta}\rangle,\]
	
	Following the convention of Definition~\ref{def:WeylGroupConvention}, the Weyl group of $\Ddeg$ is the Weyl group of $\fsl_2\oplus\fsl_2\oplus \fsl_2$, i.e.\ $W \simeq (\mathbb{Z}/2\mathbb{Z})^3$ generated by the reflections $s_1,s_2,s_3$ with $s_i(\varepsilon_i)=-\varepsilon_i$.
	
	\subsection{Affinization}\label{ssec:Affinization}
	The main character of this paper is the affinization of the Lie superalgebra $\Ddeg$.
	
	Let $z$ be a variable and let $\Bbbk[z,z^{-1}]$ be the algebra of Laurent polynomials in this variable. We first consider the loop algebra
	\begin{equation*}
		\Ddeg[z,z^{-1}]:=\Ddeg \otimes \Bbbk[z,z^{-1}].
	\end{equation*}
	
	The odd part of this superalgebra is an abelian radical and its even part is isomorphic to $(\fsl_2\oplus \fsl_2 \oplus \fsl_2)\otimes \Bbbk[z,z^{-1}]$. We denote by $\kappa(\cdot,\cdot)$ the Killing form of the semisimple Lie algebra $\fsl_2\oplus \fsl_2 \oplus \fsl_2$. Recall that it is non-degenerate. We extend this bilinear form to the whole superalgebra $\Ddeg$ by zero on the odd part, i.e.
	\begin{equation*}
		\kappa(x,y)=
		\begin{cases}
			\tr ({\rm ad}\, x\, {\rm ad}\, y),&~\text{if } x,y \in \Ddeg_{\overline 0},\\
			0,&~\text{otherwise.}
		\end{cases}
	\end{equation*}
	
	Finally, we define the affine algebra $\Daff$ as follows:
	\[\Daff := \Ddeg\otimes \Bbbk[z,z^{-1}] \oplus \Bbbk K\]
	as a vector space, where $K$ is a central even element, and
	\[[x \otimes z^a, y \otimes z^b]=[x,y]\otimes z^{a+b}+\delta_{a+b,0}\,a\,\kappa(x,y)K.\]
	Let $\Phi^{af}$ be the root system of this affine algebra. It is a disjoint union of an even and an odd part:
	\[\Phi^{af}=\Phi^{af}_{\overline 0} \sqcup \Phi^{af}_{\overline 1}.\]
	
	We have:
	\[\Phi^{af}_{\overline 0}=\{\pm 2 \varepsilon_i+ r \delta,~ r \delta \mid i=1,2,3,~ r \in \mathbb{Z}\},\]
	\[\Phi^{af}_{\overline 1}=\{\pm \varepsilon_1 \pm \varepsilon_2 \pm \varepsilon_3+ r \delta \mid r \in \mathbb{Z}\}.\]
	
	All the real roots of $\Phi^{af}$ have multiplicity one. The imaginary roots $r\delta$, $r \neq 0$, have multiplicity three: the corresponding root space is spanned by $h_1z^r, h_2z^r,h_3z^r$. This multiplicity is the source of the factor $(q;q)^{-3}$ in the character of the enveloping algebra of the radical computed in Subsection~\ref{ssec:ProjectiveModules}.
	% \Daniil{Check this.}
	% \Eugene{Strongly agree.}
	
	We define the Iwahori (affine Borel) subalgebra $\mathcal{I} \subset \Daff$ by
	\[\mathcal{I}:=\Ddeg \otimes z\Bbbk[z]\oplus\fb\oplus \Bbbk K.\]
	
	We denote by $\Phi^{af}_+$ the positive part of the root system $\Phi^{af}$ with respect to $\mathcal{I}$. That is, $\gamma \in \Phi^{af}_+$ if and only if $e_{\gamma} \in \mathcal{I}$.
	
	Define the simple affine root
	\[\alpha_0=-\varepsilon_1 - \varepsilon_2 -\varepsilon_3 + \delta\]
	and the corresponding root vector
	\[e_{\alpha_0}=v_-\otimes v_- \otimes v_- \otimes z.\]
	Note that $\varepsilon_1 + \varepsilon_2 +\varepsilon_3$ is the highest root of the finite root system. Therefore the following lemma holds.
	
	\begin{lemma}
		\[\Phi^{af}_+=\left(\bigoplus_{i=0}^3 \mathbb{Z}_{\geq 0}\alpha_i\right)\cap\Phi^{af}.\]
	\end{lemma}
	
	For any subset $J \subset \{0,1,2,3\}$ we define a parabolic subalgebra as follows. First we define the corresponding subset of the root system:
	
	\begin{definition}\label{def:parabolicRoots}
		\[\Phi^{af}_{J}:=\Phi^{af}_+\cup\left(\left(\bigoplus_{i\in J} \mathbb{Z}_{\geq 0}(-\alpha_i)\right)\cap\Phi^{af}\right).\]
	\end{definition}
	
	Then
	
	\begin{definition}\label{def:parabolicSubalgebras}
		\[\fp_J:=\fh \oplus \Bbbk K \oplus \bigoplus_{\gamma \in \Phi^{af}_{J}}\Bbbk e_{\gamma}.\]
	\end{definition}
	
	By definition we have 
	\begin{equation*}
		\fp_{\emptyset} =\mathcal{I},~\fp_{123}=\Ddeg \otimes \Bbbk[z] \oplus \Bbbk K,~ \fp_{J_1} \subset \fp_{J_2}~\text{if and only if } J_1 \subset J_2.
	\end{equation*}
	
	\begin{remark}\label{rem:ParabolicDegeneration}
		In the semisimple case a parabolic subalgebra can be defined as the subalgebra generated by the Cartan subalgebra, all positive root vectors and a chosen set of simple negative root vectors. This does not work in the degenerate case we are dealing with. Our definition, however, ensures that our parabolic subalgebras are the degenerations of the corresponding parabolic subalgebras of the semisimple Lie superalgebra $D^{(1)}(2|1,a_1,a_2,a_3)$ for generic values of the parameters.
	\end{remark}
	
	\subsection{Diagrammatic automorphisms}\label{ssec:DiagrammaticAutomorphisms}
	The goal of this subsection is to construct a family of automorphisms of the affine Lie superalgebra $\Daff$ preserving its Iwahori subalgebra $\mathcal{I}$. We call such automorphisms {\it diagrammatic} to emphasize that they are analogues of the automorphisms of semisimple Lie (super)algebras coming from automorphisms of Dynkin diagrams.
	
	The diagrammatic automorphisms we construct are indexed by the symmetric group $\mathfrak S_4$, acting naturally on the set $\{0,1,2,3\}$ of indices of the simple roots. We define the corresponding action on the root system first. For $\sigma \in \mathfrak S_4$ we set
	\[\pi_\sigma(\alpha_i):=\alpha_{\sigma(i)},\]
	and extend this to an arbitrary $\gamma=\sum_{i=0}^3a_i\alpha_i \in \Phi^{af}$ by
	\begin{equation*}\label{eq:AutomorphismOnRoot}
		\pi_\sigma(\gamma):=\sum_{i=0}^3a_i\alpha_{\sigma(i)}.
	\end{equation*}
	
	\begin{lemma}
		$\pi_\sigma(\gamma)\in \Phi^{af}$. If $\gamma \in \Phi^{af}_+$ then $\pi_\sigma(\gamma)\in \Phi^{af}_+$.
	\end{lemma}
	\begin{proof}
		Note that $\delta=\alpha_0+\alpha_1+\alpha_2+\alpha_3$, so $\pi_\sigma(\delta)=\delta$. One can rewrite
		\[\gamma=\sum_{i=1}^3(a_i-a_0)\alpha_i+a_0 \delta,\]
		and the condition $\gamma \in \Phi^{af}$ is equivalent to $\gamma^{re}:=\sum_{i=1}^3(a_i-a_0)\alpha_i \in \Phi.$
		Hence it suffices to check that $\pi_\sigma(\gamma^{re}) \in \Phi^{af}$. This follows from the description $\Phi=\{\pm \alpha_i, \pm (\alpha_0-\delta),\pm (\alpha_i +\alpha_j)\mid i,j=1,2,3\}$, which gives $\pi_\sigma(\Phi) \subset \Phi^{af}$.
		
		The second claim of the lemma is now immediate from the definition.
	\end{proof}
	
	\begin{proposition}
		The action of $\mathfrak{S}_4$ on $\Phi^{af}$ extends naturally to an action on the Lie superalgebra $\Daff$ by automorphisms. This extension satisfies
		$\pi_\sigma(e_{\alpha_i})=e_{\alpha_{\sigma(i)}}$, $i=0,1,2,3$.
	\end{proposition}
	\begin{proof}
		Assume first that $\sigma(0)=0$. Then we set
		\[\pi_\sigma(e_{\pm 2 \varepsilon_i}z^k):=e_{\pm 2 \varepsilon_{\sigma(i)}}z^k,\qquad \pi_\sigma(e_{\pm \alpha_i}z^k):=e_{\pm \alpha_{\sigma(i)}}z^k.\]
		
		This map is an automorphism by construction. Now let $\pi:=\pi_{(0,1)}$ be the transposition of $0$ and $1$. We set
		\[\pi(e_{\pm 2 \varepsilon_1}z^k):=e_{\pm 2 \varepsilon_1}z^k,~\pi(e_{\pm \alpha_i}z^k):=e_{\pm \alpha_i}z^k,~i=2,3,\]
		\[\pi(e_{2 \varepsilon_2}z^k):=e_{-2 \varepsilon_3}z^{k+1},~\pi(e_{2 \varepsilon_3}z^k):=e_{-2 \varepsilon_2}z^{k+1},\]
		\[\pi(e_{-2 \varepsilon_2}z^k):=e_{2 \varepsilon_3}z^{k-1},~\pi(e_{-2 \varepsilon_3}z^k):=e_{2 \varepsilon_2}z^{k-1},\]
		\[\pi(v_- \otimes v_+ \otimes v_+ z^k):=v_- \otimes v_- \otimes v_- z^{k+1},~\pi(v_- \otimes v_- \otimes v_- z^k):=v_- \otimes v_+ \otimes v_+ z^{k-1}.\]
		It remains to define the action on the Cartan subalgebra and on the central element and to check the compatibility with the cocycle. On the elements $h_iz^a$ with $a \neq 0$ we set $\pi_\sigma(h_iz^a):=\overline{\pi}_\sigma(h_i)z^a$, where $\overline{\pi}_\sigma$ is the linear map on $\fh$ transpose-inverse to the map induced by $\pi_\sigma$ on $\fh^*$. Explicitly, for $\pi=\pi_{(0,1)}$ the map on $\fh^*$ is $\varepsilon_1 \mapsto \varepsilon_1$, $\varepsilon_2 \mapsto -\varepsilon_3$, $\varepsilon_3 \mapsto -\varepsilon_2$, and hence
		\[\overline{\pi}(h_1)=h_1,\quad \overline{\pi}(h_2)=-h_3,\quad \overline{\pi}(h_3)=-h_2.\]
		On the finite Cartan subalgebra itself the naive formula $\pi(h):=\overline{\pi}(h)$ is {\it not} compatible with the central extension. Indeed, for $\gamma=2\varepsilon_2$ we have $\pi(e_{\gamma}z^k)=e_{-2\varepsilon_3}z^{k+1}$ and $\pi(e_{-\gamma}z^{-k})=e_{2\varepsilon_3}z^{-k-1}$, so that
		\[\pi\bigl([e_{\gamma}z^k,e_{-\gamma}z^{-k}]\bigr)=[e_{-2\varepsilon_3}z^{k+1},e_{2\varepsilon_3}z^{-k-1}]=-h_3+(k+1)\kappa(e_2,f_2)K,\]
		while $[e_\gamma z^k,e_{-\gamma}z^{-k}]=h_2+k\kappa(e_2,f_2)K$. Therefore we set $\pi_\sigma(K):=K$ and
		\begin{equation*}\label{eq:PiOnCartan}
			\pi_\sigma(h):=\overline{\pi}_\sigma(h)+\lambda_\sigma(h)K,\qquad h \in \fh,
		\end{equation*}
		where $\lambda_\sigma \in \fh^*$ is determined by $\lambda_\sigma(h_{\gamma})=n_{\gamma}\,\kappa(e_\gamma,e_{-\gamma})$. Here $\gamma$ runs over the even roots of the finite root system, $h_\gamma=[e_\gamma,e_{-\gamma}]$ and $n_\gamma \in \mathbb{Z}$ is defined by $\pi_\sigma(\gamma)=\overline{\pi}_\sigma(\gamma)+n_\gamma\delta$. For $\sigma(0)=0$ one has $n_\gamma=0$ and $\lambda_\sigma=0$, whereas for $\pi=\pi_{(0,1)}$ one gets
		\[\pi(h_1)=h_1,\quad \pi(h_2)=-h_3+\kappa(e_2,f_2)K,\quad \pi(h_3)=-h_2+\kappa(e_3,f_3)K.\]
		With these formulas a direct check on the pairs of root vectors shows that $\pi_\sigma$ preserves all the brackets, including the cocycle term. Note that the odd part of $\Ddeg$ is abelian and $\kappa$ vanishes on it, so the odd root vectors impose no condition on $\lambda_\sigma$. Finally, the relations of $\mathfrak{S}_4$ are checked directly.
	\end{proof}
	% \Daniil{Claude added some arguments for Cartan and central element, check it.}
	
	\begin{corollary}\label{cor:ParabolicStable}
		Consider the parabolic subalgebra $\fp_{ij}$ for a two-element subset $\{i,j\}\subset \{0,1,2,3\}$ and denote $\{k,l\}:=\{0,1,2,3\} \backslash \{i,j\}$. Then $\fp_{ij}$ is stable under $\pi_{(i,j)(k,l)}$.
	\end{corollary}
	
	\section{Categories of representations and Grothendieck groups}\label{sec:Categories}
	In this section we first recall the standard homological notions which we use, referring to \cite{weibel1994introduction} for the details, then the definition of the Grothendieck group and of the maps induced by functors on it. After that we introduce the categories of representations we work with and compute their Grothendieck groups in terms of characters. Finally we introduce the $\Ext$ pairing, which is the main tool of Section~\ref{sec:Orthogonality}.
	
	\subsection{Recollected facts from homological algebra}
	Let $\mathcal{O}$ be an abelian category.
	A \textbf{chain complex} \( M^\bullet \) in the abelian category \( \mathcal{O} \) is a family of objects \( \{M^n\}_{n \in \mathbb{Z}} \) in \( \mathcal{O} \) together with morphisms (differentials)
	\[
	d_M^n: M^n \longrightarrow M^{n+1}
	\]
	such that
	\[
	d_M^{n+1} \circ d_M^n = 0
	\]
	for every \( n \in \mathbb{Z} \).
	
	A \textbf{morphism} of chain complexes \( f: M^\bullet \to N^\bullet \) is a family of morphisms \( f^n: M^n \to N^n \) such that
	\[
	f^{n+1} \circ d_M^n = d_N^n \circ f^n
	\]
	for all \( n \).
	
	
	Let \( K(\mathcal{O}) \) be the \textbf{homotopy category}: objects are chain complexes, and morphisms are chain maps modulo chain homotopy.
	
	For any object $M^\bullet$ the $i$-th cohomology group is defined in the following way
	
	\[H^{i}(M)=\ker d_M^{i}/{\rm im} ~d_{M}^{i-1}.\]
	
	A chain map is a \textbf{quasi-isomorphism} if it induces isomorphisms on all homology objects.
	
	The \textbf{derived category} \( D(\mathcal{O}) \) is the localization of \( K(\mathcal{O}) \) with respect to all quasi-isomorphisms:
	\[
	D(\mathcal{O}) := K(\mathcal{O})[\mathcal{Q}^{-1}],
	\]
	where \( \mathcal{Q} \) is the class of quasi-isomorphisms.
	
	The subcategory of the derived category which consists of the elements which have right bounded representatives is denoted by $D^-(\mathcal{O})$. In a similar way $D^+(\mathcal{O})$ consists of classes of left bounded complexes, $D^b(\mathcal{O})$ of the bounded complexes.
	
	Let \(\psi: \mathcal{O} \to \mathcal{A}\) be a \textbf{right exact} functor between abelian categories, where \(\mathcal{O}\) has enough projectives.
	
	
	The \textbf{left derived functor} of \(\psi\) is the functor
	\[
	L\psi: D^-(\mathcal{O}) \longrightarrow D^-(\mathcal{A})
	\]
	defined as follows:
	
	For any complex \(X^\bullet \in D^-(\mathcal{O})\), choose a \textbf{projective resolution} 
	\[
	P^\bullet \xrightarrow{\sim} X^\bullet,
	\]
	where \(P^\bullet\) is a complex of projective objects in \(\mathcal{O}\) quasi-isomorphic to \(X^\bullet\).
	
	Then set
	\[
	L\psi(X^\bullet) := \psi(P^\bullet) \in D^-(\mathcal{A}),
	\]
	where \(\psi(P^\bullet)\) is the complex obtained by applying \(\psi\) object-wise to \(P^\bullet\).
	For a morphism in \(D^-(\mathcal{O})\), \(L\psi\) is defined by lifting it to projective resolutions. This is well-defined up to unique isomorphism.
	
	For an object \(X \in \mathcal{O}\) (viewed as a complex concentrated in degree 0), the homology of \(L\psi(X)\) gives the classical left derived functors:
	\[
	H^{-n}(L\psi(X)) \cong L^n\psi(X).
	\]
	Recall the following definition of adjoint functors.
	\begin{definition}\label{def:Adjunction}    
		An adjunction between \(\mathcal{O}\) and \(\mathcal{C}\) is given by functors 
		\(\psi: \mathcal{O} \to \mathcal{C}\) and \(\phi: \mathcal{C} \to \mathcal{O}\) 
		together with a natural bijection between Hom-sets
		\[
		\operatorname{Hom}_{\mathcal{C}}(\psi(M), N) \cong \operatorname{Hom}_{\mathcal{O}}(M, \phi(N))
		\]
		for all objects \(M \in \mathcal{O}\) and \(N \in \mathcal{C}\). 
		The bijection is required to be \textbf{natural} in both variables \(M\) and \(N\).
	\end{definition}
	
	% \Eugene{Here will be written the used facts about abelian categories, derived categories, derived functors, adjoint functors.}
	
	\subsection{Grothendieck groups}\label{ssec:GrothendieckGroups}
	Consider an abelian category $\mathcal{O}$. Recall that the Grothendieck group $K_0(\mathcal{O})$ is an abelian group generated by the symbols $[M]$, $M \in \mathcal{O}$ and for any short exact sequence 
	\begin{equation}\label{eq:ShortExactSequence}
		0 \rightarrow L \rightarrow M \rightarrow N \rightarrow 0
	\end{equation}
	there is the following relation in $K_0(\mathcal{O})$:
	\begin{equation*}
		[M]=[L]+[N].
	\end{equation*}
	In words it means that the class of an object is equal to the sum of its subobject and quotient object. We keep the letter $K$ for the homotopy category of the previous subsection and always write $K_0$ for the Grothendieck group.
	
	Recall also that the natural map $K_0(\mathcal{O}) \rightarrow K_0(D^b(\mathcal{O}))$, $[M]\mapsto [M]$, is an isomorphism, the inverse being given by the Euler characteristic $[M^\bullet]\mapsto \sum_i(-1)^i[H^i(M^\bullet)]$. We use this identification without further mention.
	
	Now let $\eta$ be an exact functor $\mathcal{O} \rightarrow \mathcal{C}$. For any exact sequence \eqref{eq:ShortExactSequence} it gives the short exact sequence
	\[
	0 \rightarrow \eta(L) \rightarrow \eta(M) \rightarrow \eta(N) \rightarrow 0.\]
	Thus 
	\[[\eta(M)]=[\eta(L)]+[\eta(N)].\]
	
	Therefore each exact functor naturally defines the map on the Grothendieck groups $[\eta]:K_0(\mathcal{O})\rightarrow K_0(\mathcal{C})$, $[M] \mapsto [\eta(M)]$.
	
	More generally, assume that $\eta$ is right exact functor of finite homological degree $k$, i.e.\ $\textbf{L}^{(k+1)}\eta=0$. Then for any short exact sequence \eqref{eq:ShortExactSequence} we get
	\[0\rightarrow  \textbf{L}^{(k)}\eta(L) \rightarrow \textbf{L}^{(k)}\eta(M) \rightarrow \textbf{L}^{(k)}\eta(N) \rightarrow  \dots\rightarrow \textbf{L}^{(1)}\eta(L) \rightarrow \textbf{L}^{(1)}\eta(M) \rightarrow \textbf{L}^{(1)}\eta(N)  \rightarrow \eta(L) \rightarrow \eta(M) \rightarrow \eta(N) \rightarrow 0.\]
	
	Then the following relation holds in the Grothendieck group:
	\[\sum_{i=0}^k (-1)^i[\textbf{L}^{(i)}\eta(M)]=\sum_{i=0}^k (-1)^i [\textbf{L}^{(i)}\eta(L)]+\sum_{i=0}^k (-1)^i [\textbf{L}^{(i)}\eta(N)].\]
	Therefore each right exact functor naturally defines the map on the Grothendieck groups
	\begin{equation*}
		[\eta]:K_0(\mathcal{O}) \rightarrow K_0(\mathcal{C}),\qquad [M] \mapsto \sum_{i=0}^k (-1)^i[\textbf{L}^{(i)}\eta(M)].
	\end{equation*}
	
	In the same way the left exact functor acts naturally on the Grothendieck group.
	
	\subsection{Categories of representations}\label{ssec:CategoriesOfRepresentations}
	We now introduce the categories of modules we work with. Let $\fg$ be a Lie superalgebra with a fixed decomposition \eqref{eq:ReductivePart}, let $\fb$ be a Borel subalgebra of $\fg^{red}$ and let $\fh \subset \fb$ be the corresponding Cartan subalgebra.
	
	\begin{definition}\label{def:CategoryO}
		We denote by $\mathcal{O}(\fg)$ the category of $\fg$-modules integrable with respect to $\fg^{red}$ such that the weights of these modules with respect to $\fh$ are bounded from the above. In a similar way we denote by $\mathcal{O}^\vee(\fg)$ the category of $\fg$-modules integrable with respect to $\fg^{red}$ such that the weights of these modules with respect to $\fh$ are bounded from the below.
	\end{definition}
	
	In particular, finite dimensional $\fg$-modules are contained in $\mathcal{O}(\fg) \cap \mathcal{O}^\vee(\fg)$.
	
	Let $R$ be the radical of $\fg$, i.e.\ $\fg \simeq \fg^{red} \ltimes R$. Then $R$ acts nilpotently on each object $M \in \mathcal{O}(\fg)$. Thus it annihilates each irreducible object. Therefore 
	\[K_0(\mathcal{O}(\fg))\simeq K_0(\mathcal{O}(\fg^{red})).\]
	
	Therefore $K_0(\mathcal{O}(\fg))$ is a completion of the free abelian group on the classes of the irreducible modules. The modules occurring in this paper are in general not of finite length, but their weight components in each fixed $z$-degree are, so that their classes are infinite sums of classes of irreducibles with only finitely many terms in each degree in $z$.	
	
	For a finite dimensional $\fg$-module $V$ the tensor product functor $\bullet \otimes V:\mathcal{O}(\fg) \rightarrow \mathcal{O}(\fg)$ is an exact endofunctor. 
	
	All the functors we consider relate a pair of Lie superalgebras $\fg_1 \subset \fg_2$. Throughout the paper such a pair is assumed to satisfy the following assumption.
	
	\begin{assumption}\label{ass}
		Let $\fb_2$ be a Borel subalgebra of $\fg_2^{red}$, let $\fb_2^-$ be the opposite Borel subalgebra and let $\fh \subset \fb_2$ be the corresponding Cartan subalgebra. We assume that 
		\begin{itemize}
			\item $\fb_2 \subset \fg_1$,
			\item $\fg_{1}^{\overline 0}+ \fb_2^-=\fg_{2}^{ \overline 0}$,
            \item $\dim(\fg_2/\fg_1) < \infty$.
		\end{itemize} 
	\end{assumption}
	
	Finally we introduce the graded version of these categories in the case which is of interest for us. Recall that the Lie superalgebras $\fp_J$ have various natural gradings. First of all they have a $\mathbb{Z}$-grading such that the even part of $\fp_J$ is contained in the zero component and the odd part of $\fp_J$ is contained in the odd component. Next this superalgebra is graded with respect to the $z$-degree. Finally it is graded by the eigenvalue of the adjoint action of the Cartan subalgebra.
	
	\begin{definition}\label{def:GradedCategory}
		We denote by $\mathcal{O}^{gr}(\fp_J)$ the category of integrable representations of $\fp_J$ graded with respect to all these gradings.
	\end{definition}
	
	For $M \in \mathcal{O}^{gr}(\fp_J)$ and $l,r \in \mathbb{Z}$ we denote by $M\sh{l}{r}$ the module $M$ with the $z$-grading shifted by $l$ and the parity grading shifted by $r$. The defining property is
	\begin{equation}\label{eq:ShiftDefinition}
		\ch\bigl(M\sh{l}{r}\bigr)=q^l(-t)^r\ch(M),
	\end{equation}
	where $\ch$ is the character defined in Subsection~\ref{ssec:Characters}. The sign is explained in Remark~\ref{rem:SignConvention}.
	The homological shift of a complex is denoted, as always, by $M^\bullet[n]$, with the convention $\bigl(M^\bullet[n]\bigr)^i=M^{i+n}$. In particular a left derived functor takes values in non-positive cohomological degrees and a right derived one in non-negative degrees. Thus the angle brackets always refer to the two internal gradings and the square brackets always refer to the homological grading. In the Grothendieck group
	\[\bigl[M\sh{l}{r}\bigr]=q^l(-t)^r[M],\qquad \bigl[M^\bullet[1]\bigr]=-[M^\bullet].\]
	
	\subsection{Characters}\label{ssec:Characters}
	Let us now compute the Grothendieck groups of the categories $\mathcal{O}^{gr}(\fp_J)$ in terms of characters.
	
	%Consider the Lie superalgebra $\Ddeg\otimes \Bbbk[z] \simeq \fp_{123}$. Then 
	%\[\fp_{123}^{red}=\fsl_1 \oplus \fsl_2 \oplus \fsl_2.\]
	Let $M \in \mathcal{O}^{gr}(\fp_J)$. $M$ is integrable over $\fp_J^{red}$. Therefore the Cartan subalgebra $\fh= h_1 \oplus h_2 \oplus h_3$ acts semisimply on $M$. Thus $M$ has the following decomposition as a vector space:
	\begin{equation*}
		M:=\bigoplus_{\lambda \in \mathbb{Z}^3,k,r\in \mathbb{Z}}M_{\lambda,k,r},
	\end{equation*}
	such that $M_{\lambda,k,r}$ lies in the $k$-th $z$-graded component, the $r$-th component coming from the parity grading and for $m \in M_{\lambda,k,r}$
	\[h_i m:=\lambda_i m.\]
	
	Then the character of the module $M$ is defined to be the graded superdimension of $M$, that is, the variable recording the odd grading is $-t$:
	\begin{equation}\label{eq:CharacterDefinition}
		\ch(M):=\sum_{\lambda \in \mathbb{Z}^3,k,r\in \mathbb{Z}} \dim(M_{\lambda,k,r})X_1^{\lambda_1}X_2^{\lambda_2}X_3^{\lambda_3}q^k(-t)^r.
	\end{equation}
	
	\begin{remark}\label{rem:SignConvention}
		The sign in \eqref{eq:CharacterDefinition} is a normalization of the variable $t$ and nothing more: replacing $t$ by $-t$ turns $\ch$ into the naive graded dimension. We use this normalization because the $\mathbb{Z}$-grading by $r$ refines the parity, the component $M_{\lambda,k,r}$ is even or odd according to the parity of $r$, so that $\ch$ becomes the supercharacter, and because with it the exterior algebra on an odd one dimensional space of weight $\beta$ has character $1-tX^{\beta}$. This is the form in which all the products below appear, see \eqref{eq:FirstInductionCharacter}, the character of $\U(R)$ in Subsection~\ref{ssec:ProjectiveModules} and the numerators \eqref{eq:Cab} of the operators of \cite{ArthamonovShakirov-2019}.
	\end{remark}
	
	
	Note that for any short exact sequence \eqref{eq:ShortExactSequence} one has
	\[\ch(M)=\ch(L)+\ch(N).\]
	
	Therefore there is the natural map of the abelian groups:
	\begin{equation*}
		\iota_{J}:K_0(\mathcal{O}^{gr}(\fp_{J})) \rightarrow \mathbb{Z}[X_1^{\pm 1},X_2^{\pm 1},X_3^{\pm 1},t^{\pm 1}]((q^{\frac{1}{2}})).
	\end{equation*}
	
	Consider now the case $J=\{1,2,3\}$.
	Up to the shifts $\sh{l}{r}$ the irreducible modules over this algebra have the form $V_{123}(\mu_1,\mu_2,\mu_3):=V_{\mu_1 \varepsilon_1} \otimes V_{\mu_2 \varepsilon_2} \otimes V_{\mu_3 \varepsilon_3}$, where $V_{\mu_i \varepsilon_i}$ is the irreducible representation of the $i$-th copy of the Lie algebra $\fsl_2$, $\mu_i\geq 0$. The general irreducible object is $V_{123}(\mu)\sh{l}{r}$. Therefore
	\begin{equation}\label{eq:character3Irreducible}
		\ch(V_{123}(\mu_1,\mu_2,\mu_3))= \prod_{i=1}^3 \frac{X_i^{\mu_i}-X_i^{-\mu_i-2}}{1-X_i^{-2}}
	\end{equation} 
	Recall the reflections $s_i, i=1,2,3$, acting in the following way
	\[s_1(X_1^{\mu_1}X_2^{\mu_2}X_3^{\mu_3}q^kt^r)=X_1^{-\mu_1}X_2^{\mu_2}X_3^{\mu_3}q^kt^r.\]
	Then the following lemma is clear.
	\begin{lemma}
		The Grothendieck group of the Lie superalgebra $\fp_{123}$ is isomorphic to
		\[\iota_{123}(K_0(\mathcal{O}^{gr}(\fp_{123})))=\mathbb{Z}[X_1^{\pm 1},X_2^{\pm 1},X_3^{\pm 1},t^{\pm 1}]^{s_1,s_2,s_3}((q^{\frac{1}{2}})).\]
	\end{lemma}
	
	In the same way consider the case $J=\{1,2\}$. In this case
	\[\fp_{12}^{red}=\langle h_1,h_2 \rangle \oplus \fsl_2^{(3)},\]
	where $\fsl_2^{(3)}$ is the third copy of $\fsl_2$.
	Up to the shifts $\sh{l}{r}$ the irreducible modules over this algebra have the form $V_{12}(\mu_1,\mu_2,\mu_3)=L_{\mu_1 \varepsilon_1}\otimes L_{\mu_2 \varepsilon_2}\otimes V_{\mu_3 \varepsilon_3}$, where $L_{\mu_i \varepsilon_i}$ is the one-dimensional representation of the Lie algebra $\langle h_i \rangle$, $\mu_1,\mu_2 \in \mathbb{Z}$, $\mu_3\geq 0$. Therefore
	\begin{equation*}
		\ch(V_{12}(\mu_1,\mu_2,\mu_3))= X_1^{\mu_1}X_2^{\mu_2} \frac{X_3^{\mu_3}-X_3^{-\mu_3-2}}{1-X_3^{-2}}
	\end{equation*} 
	
	The same computations hold for $\fp_{13}$ and $\fp_{23}$. Thus we have
	\begin{lemma}
		Let $\{i,k\}=\{1,2,3\}\backslash \{j\}$. The Grothendieck group of the Lie superalgebra $\fp_{ik}$ is isomorphic to
		\[\iota_{ik}(K_0(\mathcal{O}^{gr}(\fp_{ik})))=\mathbb{Z}[X_1^{\pm 1},X_2^{\pm 1},X_3^{\pm 1},t^{\pm 1}]^{s_j}((q^{\frac{1}{2}})).\]
	\end{lemma}
	
	% 	Consider the category $\mathcal O({\fh})$ of finite-dimensional $\fh$-modules graded with respect to weight grading and the additional grading. Consider $M \in \mathcal O(\fh)$. Let $m_1, \dots, m_l$ be a homogeneous basis of $M$. Let $a_1^i\varepsilon_1+a_2^i\varepsilon_2+a_3^i\varepsilon_3+b^i\zeta$ be a weight of an element $m_i$. We denote by the character of $M$ the sum of formal exponents of these weights. More precisely, denote $X_i:=X^{\varepsilon_i}$, $t:=X^{\zeta}$. Then
	% 	\[ch(M):=\sum_{j=1}^l X_1^{a_1^i}X_2^{a_2^i}X_3^{a_3^i}t^{b^i}.\]
	
	\begin{example} We list the characters of the algebras we deal with.
		\[\ch(\mathcal{I})=3+X_1^2+X_2^2+X_3^2-tX_1^{-1}X_2X_3-tX_1X_2^{-1}X_3-tX_1X_2X_3^{-1}-tX_1X_2X_3+\frac{q}{1-q}\ch \Ddeg,\]
		\begin{equation*}
			\label{eq:ParabolicAlgebraCharacter}
			\ch(\fp_{12})=\ch(\mathcal{I})-tX_1^{-1}X_2X_3^{-1}-tX_1X_2^{-1}X_3^{-1}+X_3^{-2},
		\end{equation*}
		\begin{equation*} \label{eq:CurrentAlgebraCharacter}   
			\ch(\fp_{123})=\frac{1}{1-q}\ch \Ddeg.
		\end{equation*}
		
		Here $\ch \Ddeg$ is the character of the finite dimensional superalgebra and
		\[\ch \Ddeg=\sum_{i=1}^3 (X_i^{2}+1+X_i^{-2})-t(X_1+X_1^{-1})(X_2+X_2^{-1})(X_3+X_3^{-1}).\]
	\end{example}
	
	\subsection{Graded \texorpdfstring{$\Hom$}{Hom} and \texorpdfstring{$\Ext$}{Ext} pairing}\label{ssec:ExtPairing}
	In this subsection we introduce the $\Ext$ pairing on the Grothendieck group of $\mathcal{O}^{gr}(\fp_J)$. It is used in Section~\ref{sec:Orthogonality}, where we compute it explicitly for $J=\{1,2,3\}$.
	
	Let $M, N \in \mathcal{O}^{gr}(\fp_J)$ be two modules. Recall the grading shifts $M\sh{l}{r}$ introduced in \eqref{eq:ShiftDefinition}. We have the natural isomorphisms
	\[\Hom\bigl(M ,N\sh{l}{r}\bigr)\simeq \Hom\bigl(M\sh{-l}{-r} ,N\bigr).\]
	
	We consider the graded sums of these vector spaces
	\[\Hom_{gr}(M ,N):=\bigoplus_{l,r \in \mathbb{Z}}\Hom\bigl(M ,N\sh{l}{r}\bigr),\qquad
	\Ext^{i}_{gr}(M ,N):=\bigoplus_{l,r \in \mathbb{Z}}\Ext^{i}\bigl(M ,N\sh{l}{r}\bigr),\]
	regarded as $\mathbb{Z}^2$-graded vector spaces, the summand indexed by $(l,r)$ sitting in degree $(l,r)$.
	
	For a $\mathbb{Z}^2$-graded vector space $V=\bigoplus_{l,r}V_{l,r}$ we write
	\begin{equation*}\label{eq:GradedSuperdimension}
		\dim_{q,t}(V):=\sum_{l,r \in \mathbb{Z}}\dim(V_{l,r})\,q^l(-t)^r,
	\end{equation*}
	with the same normalization of $t$ as in \eqref{eq:CharacterDefinition}.
	
	\begin{definition}
		\[(M,N)_{\Ext}:=\sum_{i=0}^{\infty}(-1)^i\dim_{q,t}(\Ext^{i}_{gr}(M ,N)) \in \mathbb{Z}\left[t^{\pm 1}\right]((q^{\frac{1}{2}})).\]
	\end{definition}
	
	Note that this sum in general may be not well defined. However the radical of our algebras have positive $q,t$ degrees, therefore such a sum is well defined formal series.
	
	Consider short exact sequences
	\[0 \rightarrow M_1 \rightarrow M \rightarrow M_2 \rightarrow 0,\]
	\[0 \rightarrow N_1 \rightarrow N \rightarrow N_2 \rightarrow 0.\]
	
	Then the $\Ext$ long exact sequence gives the following equalities:
	\begin{equation*}\label{eq:ExtPairingShortDecomposition}
		(M,N)_{\Ext}=(M_1,N)_{\Ext}+(M_2,N)_{\Ext}=(M,N_1)_{\Ext}+(M,N_2)_{\Ext}.
	\end{equation*}
    See for example \cite{weibel1994introduction}.
	
	Moreover by definition the $\Ext$ pairing is $q,t$-skew linear:
	\begin{equation*}
		\bigl(M\sh{l}{r},N\bigr)_{\Ext}=(M,N)_{\Ext}\,q^l(-t)^r=\bigl(M,N\sh{-l}{-r}\bigr)_{\Ext}.
	\end{equation*}
	
	Therefore the value of the $\Ext$ pairing depends only on the class of the module in the Grothendieck group.
	
	We first check that the category we work with has enough projectives and describe them.
	
	\begin{lemma}\label{lem:Projectivity}
		For any irreducible $\fp_J^{red}$-module $V$ the module
		\[P(V):=\U(\fp_J)\otimes_{\U(\fp_J^{red})}V\]
		is projective in $\mathcal{O}^{gr}(\fp_J)$. The category $\mathcal{O}^{gr}(\fp_J)$ has enough projectives, and $P_J(\mu):=P(V_J(\mu))$ is the projective cover of $V_J(\mu)$.
	\end{lemma}
	\begin{proof}
		By the tensor-hom adjunction, for any $M \in \mathcal{O}^{gr}(\fp_J)$ we have a natural isomorphism
		\[\Hom_{\fp_J}(P(V),M)\simeq \Hom_{\fp_J^{red}}\bigl(V,\Res^{\fp_J}_{\fp_J^{red}}M\bigr).\]
		The algebra $\fp_J^{red}$ is reductive and $M$ is integrable over it, hence $\Res^{\fp_J}_{\fp_J^{red}}M$ is a semisimple $\fp_J^{red}$-module and the functor $\Hom_{\fp_J^{red}}(V,\bullet)$ is exact on the essential image of the restriction. Therefore $\Hom_{\fp_J}(P(V),\bullet)$ is exact, i.e.\ $P(V)$ is projective.
		
		As a $\fp_J^{red}$-module $P(V)\simeq \U(R)\otimes V$, where $R$ is the radical of $\fp_J$. Since $R$ acts nilpotently, $P(V)$ has a unique maximal submodule, namely $R\cdot P(V)$, and $P(V)/R\cdot P(V)\simeq V$. Hence $P(V)$ is the projective cover of the irreducible module with the same highest weight.
		
		Finally, every $M \in \mathcal{O}^{gr}(\fp_J)$ is a direct sum of irreducible submodules as $\fp_J^{red}$ module, so there is a surjection onto $M$ from a direct sum of modules of the form $P_J(\mu)\sh{l}{r}$. Thus the category has enough projectives.
	\end{proof}
	
	For an irreducible module $V_{J}(\mu)$ let $P_{J}(\mu)$ be its projective cover, as in Lemma~\ref{lem:Projectivity}. Projective modules are $\Ext$-acyclic, i.e.\
	\[\Ext^{i}(P_{J}(\mu) ,N)=\{0\},~ i \geq 1.\]
	Moreover
	\[\dim \Hom\bigl(P_{J}(\mu)\sh{l}{r},V_{J}(\mu')\sh{l'}{r'} \bigr)=\delta_{\mu,\mu'}\delta_{l,l'}\delta_{r,r'}.\]
	
	Therefore one gets
	\begin{equation*}
		\bigl(P_{J}(\mu)\sh{l}{r},V_{J}(\mu')\sh{l'}{r'}\bigr)_{\Ext} =q^{l-l'}(-t)^{r-r'}\delta_{\mu,\mu'}.
	\end{equation*}
	
	Let $(-,-): K_0( \mathcal{O}^{gr}(\fp_{J}))\times K_0( \mathcal{O}^{gr}(\fp_{J})) \rightarrow \mathbb{Z}[t^{\pm 1}]((q^{\pm \frac{1}{2}}))$ be the pairing on the Grothendieck group induced by the $\Ext$ pairing:
	\begin{equation*}
		([M],[N]):=(M,N)_{\Ext}.
	\end{equation*}
	
	Then this is a skew $\mathbb{Z}[q^{\pm \frac{1}{2}},t^{\pm 1}]$-bilinear form, i.e.\
	\[([M]q^lt^r,[N])=([M],[N])q^lt^r=([M],[N]q^{-l}t^{-r}).\]
	
	The classes of projective modules form a dual basis for the classes of irreducible modules, i.e.\
	\begin{equation*}
		([P_{J}(\mu)],[V_{J}(\mu')])=\delta_{\mu,\mu'}.
	\end{equation*}
	
	This pairing naturally extends to the derived categories. For $M^\bullet, N^\bullet \in D^b(\mathcal{O}^{gr}(\fp_{J}))$ one can define 
	
	\begin{definition}
		\[(M^\bullet,N^\bullet )_{\Ext}:=\sum_{i=-\infty}^{\infty}(-1)^i\dim_{q,t}\bigl(\Hom_{gr}(M^\bullet ,N^\bullet[i])\bigr),\]
		where $[i]$ is the homological shift. 
	\end{definition}
	
	If the modules $M,N$ are identified with the complexes concentrated in zero degree, then both definitions of $\Ext$-pairing are equivalent.
	
	\section{Functors on the categories of representations}\label{sec:Functors} 
	The goal of this section is to study the functors between categories of representations of the parabolic subalgebras of $\Daff$. The main results of this paper mean that these functors behave in the same way as the functors between certain parabolic subalgebras of the Lie superalgebra $\fg \otimes \Bbbk[\xi]$, where $\xi$ is an odd variable and $\fg$ is a Kac--Moody Lie algebra, see \cite{KKhM}. This similarity gives us a way to see why the algebra of the decategorifications of these functors behaves similarly to the classical Double Affine Hecke Algebra \cite{Cherednik-2005}. For the general theory of Zuckerman functors, see for example \cite{Vogan1981}.
	
	\subsection{Restriction and induction functors}
	Consider a pair of Lie superalgebras $\fg_1 \subset \fg_2$ and let $\rho$ be a natural inclusion. We first recall the three standard functors relating the categories of modules over $\fg_1$ and over $\fg_2$, and then their integrable versions, which are the ones we actually use.
	
	\begin{definition}
		The restriction functor $\Res_{\fg_1}^{\fg_2}$ sends a $\fg_2$-module $M$ to the same vector space regarded as a $\fg_1$-module, and a morphism of $\fg_2$-modules to the same map of vector spaces.
	\end{definition}
	
	This functor is exact and preserves finite-dimensionality.
	
	\begin{definition}
		For a \(\fg_1\)-module \(V\) the \textbf{induced} and the \textbf{coinduced} modules are
		\[
		\Ind_{\fg_1}^{\fg_2}(V) = \U(\fg_2) \otimes_{\U(\fg_1)} V,\qquad
		\coInd_{\fg_1}^{\fg_2}(V) = \Hom_{\U(\fg_1)}(\U(\fg_2),V).
		\]
	\end{definition}
	
	Let $\fg_2^{red}$ be the reductive part of $\fg_2$ in the sense of \eqref{eq:ReductivePart}. Recall that it is contained in the even subalgebra $\fg_{2}^{\overline 0}$. Let $M$ be a $\fg_2$-module.
	\begin{definition}
		We denote by $M_{int}$ the maximal $\fg_2^{red}$ integrable quotient of $M$, and by $M^{int}$ the maximal $\fg_2^{red}$ integrable submodule of $M$.
	\end{definition}
	
	Note that at this moment we have not discussed the existence of $M_{int}$ and $M^{int}$.
	
	Next we define integrable induction and coinduction functors.
	
	\begin{definition}
		The integrable induction and coinduction functors are
		\[
		\widetilde \Ind_{\fg_1}^{\fg_{2}}(V) = \bigl(\U(\fg_2) \otimes_{\U(\fg_1)} V\bigr)_{int},\qquad
		\widetilde \coInd_{\fg_1}^{\fg_{2}}(V) = \bigl(\Hom_{\U(\fg_1)}(\U(\fg_2),V)\bigr)^{int}.
		\]
	\end{definition}
	
	Recall that the categories $\mathcal{O}(\fg_i)$, $\mathcal{O}^\vee(\fg_i)$, $i=1,2$, were introduced in Definition~\ref{def:CategoryO} and that the pair $\fg_1 \subset \fg_2$ is assumed to satisfy Assumption~\ref{ass}.
	
	\begin{lemma}
		\begin{itemize}
			\item $\Res_{\fg_1}^{\fg_2}(\mathcal{O}(\fg_2))\subset \mathcal{O}(\fg_1)$, $\Res_{\fg_1}^{\fg_2}(\mathcal{O}^\vee(\fg_2))\subset \mathcal{O}^\vee(\fg_1)$,
			\item $\widetilde \Ind_{\fg_1}^{\fg_{2}}\mathcal{O}(\fg_1) \subset \mathcal{O}(\fg_2)$,
			\item $\widetilde \coInd_{\fg_1}^{\fg_{2}}\mathcal{O}^\vee(\fg_1) \subset \mathcal{O}^\vee(\fg_2)$.
		\end{itemize}
	\end{lemma}
	\begin{proof}
		The first claim is obvious. For the second claim recall the Cartan subalgebra $\fh \subset \fg_1$ of Assumption~\ref{ass}. Since $\fh$ acts semisimply on $\fg_2$ by
		the adjoint action and preserves $\fg_1$, we may write $\fg_2=\fg_1\oplus \fn$, where $\fn$ is a sum of root spaces of $\fg_2$. By Assumption~\ref{ass} it is finite
		dimensional. Then for any $M \in \mathcal{O}(\fg_1)$ we have an isomorphism of $\fh$-modules
		\[\Ind_{\fg_1}^{\fg_{2}}(M)\simeq S(\fn)\otimes M,\]
		where $S(\fn)$ is a symmetric algebra of a superspace. In particular, $S(\fn)=S(\fn^{\overline 0})\otimes S(\fn^{\overline 1})$. However $S(\fn^{\overline 1})$ is finite-dimensional and $\fn^{\overline 0}$ contains only negative weights. Therefore the boundedness of weights holds and the maximal integrable quotient of this module is contained in $\mathcal{O}(\fg_2)$.
		
		In the same way for $M \in \mathcal{O}^\vee (\fg_1)$:
		\[\coInd_{\fg_1}^{\fg_{2}}(M)\simeq S(\fn)^\vee\otimes M,\]
		and all the weights of $S(\fn^{\overline 0})^\vee$ are positive. Therefore the boundedness of weights holds and the maximal integrable submodule of this module is contained in $\mathcal{O}^\vee(\fg_2)$.
	\end{proof}
	
	\begin{lemma}
		$\widetilde\Ind_{\fg_1}^{\fg_2}:\mathcal{O}(\fg_1)\rightarrow \mathcal{O}(\fg_2)$ is well defined, as well as $\widetilde\coInd_{\fg_1}^{\fg_2}:\mathcal{O}^\vee(\fg_1)\rightarrow \mathcal{O}^\vee(\fg_2)$.
	\end{lemma}
	\begin{proof}
		We prove the first claim, the proof of the second claim is similar. We need to check that $\Ind_{\fg_1}^{\fg_2}(M)$ contains a unique maximal submodule $N$ such that $\Ind_{\fg_1}^{\fg_2}(M)/N$ is integrable.
		
		Note that the weights of $\Ind_{\fg_1}^{\fg_2}(M)$ are bounded from above. Let $S \subset \Ind_{\fg_1}^{\fg_2}(M)$ be a subspace spanned by all the elements of the weights $\lambda$ such that for some $\sigma$ from the Weyl group $\sigma(\lambda)$ is greater than any weight of $\Ind_{\fg_1}^{\fg_2}(M)$. Therefore this space is contained in the kernel of any surjection $\Ind_{\fg_1}^{\fg_2}(M)\rightarrow V$, where $V$ is integrable. We denote by $N$ the submodule of $\Ind_{\fg_1}^{\fg_2}(M)$ generated by $S$, i.e.\ $N:=\U(\fg_2)S$. We claim that $\Ind_{\fg_1}^{\fg_2}(M)/N$ is integrable. Indeed its weights are contained in the convex hull of the Weyl group image of the weights of $M$. Therefore as a module over the reductive Lie algebra $\fg_2^{red}$ it is integrable.
	\end{proof}
	
	The next proposition contains a list of properties of these functors.
	
	\begin{proposition}\label{prop:Adjoint}
		\begin{itemize}
			\item $\widetilde\Ind_{\fg_1}^{\fg_2}:\mathcal{O}(\fg_1)\rightarrow \mathcal{O}(\fg_2)$ is left adjoint to $\Res_{\fg_1}^{\fg_2}:\mathcal{O}(\fg_2)\rightarrow \mathcal{O}(\fg_1)$, in particular, $\widetilde\Ind_{\fg_1}^{\fg_2}$ is right exact,
			\item $\widetilde\coInd_{\fg_1}^{\fg_2}:\mathcal{O}^\vee(\fg_1)\rightarrow \mathcal{O}^\vee(\fg_2)$ is right adjoint to $\Res_{\fg_1}^{\fg_2}:\mathcal{O}^\vee(\fg_2)\rightarrow \mathcal{O}^\vee(\fg_1)$, in particular, $\widetilde\coInd_{\fg_1}^{\fg_2}$ is left exact.
		\end{itemize}    
	\end{proposition}
	\begin{proof}
		Let $M \in \mathcal{O}(\fg_1)$ and $N \in \mathcal{O}(\fg_2)$. The usual tensor-hom adjunction gives a natural isomorphism
		\[\Hom_{\fg_2}\bigl(\Ind_{\fg_1}^{\fg_2}M,N\bigr)\simeq \Hom_{\fg_1}\bigl(M,\Res_{\fg_1}^{\fg_2}N\bigr).\]
		Since $N$ is integrable over $\fg_2^{red}$, every morphism $\Ind_{\fg_1}^{\fg_2}M \rightarrow N$ vanishes on the kernel of the projection onto the maximal integrable quotient, i.e.\ factors uniquely through $\widetilde\Ind_{\fg_1}^{\fg_2}M$. Hence
		\[\Hom_{\mathcal{O}(\fg_2)}\bigl(\widetilde\Ind_{\fg_1}^{\fg_2}M,N\bigr)\simeq \Hom_{\fg_2}\bigl(\Ind_{\fg_1}^{\fg_2}M,N\bigr)\simeq\Hom_{\mathcal{O}(\fg_1)}\bigl(M,\Res_{\fg_1}^{\fg_2}N\bigr),\]
		and all the isomorphisms are natural in $M$ and $N$. This is the required adjunction in the sense of Definition~\ref{def:Adjunction}. A left adjoint functor is right exact.
		
		The second claim is proved in the same way, starting from the isomorphism
		\[\Hom_{\fg_2}\bigl(N,\coInd_{\fg_1}^{\fg_2}M\bigr)\simeq \Hom_{\fg_1}\bigl(\Res_{\fg_1}^{\fg_2}N,M\bigr)\]
		and using that the image of an integrable module is contained in the maximal integrable submodule, so that every morphism $N \rightarrow \coInd_{\fg_1}^{\fg_2}M$ factors uniquely through $\widetilde\coInd_{\fg_1}^{\fg_2}M$. A right adjoint functor is left exact.
	\end{proof}
	
	This proposition allows one to define the derived functors
	\[\Lind_{\fg_1}^{\fg_2}:D^-(\mathcal{O}(\fg_1))\rightarrow D^-(\mathcal{O}(\fg_2)),\]
	\[\Rcoind_{\fg_1}^{\fg_2}:D^+(\mathcal{O}^\vee(\fg_1))\rightarrow D^+(\mathcal{O}^\vee(\fg_2)).\]
	
	Note that the derived functors naturally act on the Grothendieck groups. We will denote these operators by
	\[[\widetilde{\Ind}_{\fg_1}^{\fg_2}]:K_0(\mathcal{O}(\fg_1))\rightarrow K_0(\mathcal{O}(\fg_2)).\]
	
	We need the following property of induction and restriction functors.
	\begin{lemma}
		The (derived) induction and coinduction commute with the tensor product by an integrable module. More precisely
		let $B$ be a finite dimensional $\fg_2$ module. Then
		\[\widetilde\Ind_{\fg_1}^{\fg_2}(\bullet \otimes \Res_{\fg_1}^{\fg_2}B)\simeq \widetilde\Ind_{\fg_1}^{\fg_2}(\bullet)\otimes B.\]
	\end{lemma}
	\begin{proof}
		For any $K \in \mathcal{O}(\fg_2)$, $A \in \mathcal{O}(\fg_1)$ we have the natural isomorphisms:
		\begin{multline*}
			\Hom_{\mathcal{O}(\fg_2)}(\widetilde \Ind_{\fg_1}^{\fg_2}(A\otimes \Res_{\fg_1}^{\fg_2}B),K)\simeq \Hom_{\mathcal{O}(\fg_1)}(A\otimes \Res_{\fg_1}^{\fg_2}B, \Res_{\fg_1}^{\fg_2}K)\simeq\Hom_{\mathcal{O}(\fg_1)}(A, \Res_{\fg_1}^{\fg_2}K \otimes (\Res_{\fg_1}^{\fg_2}B)^*)\simeq\\
			\simeq \Hom_{\mathcal{O}(\fg_1)}(A, \Res_{\fg_1}^{\fg_2}(K \otimes  B^*))\simeq
			\Hom_{\mathcal{O}(\fg_2)}(\widetilde \Ind_{\fg_1}^{\fg_2}A,K \otimes  B^*)\simeq     \Hom_{\mathcal{O}(\fg_2)}(\widetilde \Ind_{\fg_1}^{\fg_2}A\otimes  B,K ).
		\end{multline*}
		Therefore we get the needed isomorphism of functors. The same argument, starting from the adjunction of Proposition~\ref{prop:Adjoint} for the coinduction, gives
		$\widetilde\coInd_{\fg_1}^{\fg_2}(\bullet \otimes \Res_{\fg_1}^{\fg_2}B)\simeq \widetilde\coInd_{\fg_1}^{\fg_2}(\bullet)\otimes B$, and both isomorphisms extend to
		the derived functors.
	\end{proof}
	
	\begin{corollary}\label{cor:CommutationMultiplicationSymmetric}
		\[[\widetilde{\Ind}_{\fg_1}^{\fg_2}]([A \otimes \Res_{\fg_1}^{\fg_2}B])=[\widetilde{\Ind}_{\fg_1}^{\fg_2}](A)\cdot [B].\]   
		In words, the operators on the Grothendieck groups commute with the multiplication by a restricted module.
	\end{corollary}
	
	
	% The following endofunctors of the categories $\mathcal{O}(\fg_1)$, $\mathcal{O}(\fg_1)^\vee$ are called {\it Zuckerman functors}.
	
	% 	\begin{definition}\label{def:GeneralZuckermann}
		% 		\[\mathcal{Z}_{\fg_1}^{\fg_2}:=\Res_{\fg_1}^{\fg_2}\widetilde \Ind_{\fg_1}^{\fg_2};\]
		% 		\[\widehat{\mathcal{Z}}_{\fg_1}^{\fg_2}:=\Res_{\fg_1}^{\fg_2}\widetilde \coInd_{\fg_1}^{\fg_2}.\]
		% 	\end{definition}
	
	% In this paper we will need the endofunctors of $\mathcal{O}(\fg_2)$, $\mathcal{O}(\fg_2)^\vee$ 
	
	
	% 	\begin{definition}\label{def:GeneralFZuckermann}
		% 		\[\mathcal{F}_{\fg_2}^{\fg_1}:=\widetilde \Ind_{\fg_1}^{\fg_2}\Res_{\fg_1}^{\fg_2};\]
		% 		\[\widehat{\mathcal{F}}_{\fg_1}^{\fg_2}:=\widetilde \coInd_{\fg_1}^{\fg_2}\Res_{\fg_1}^{\fg_2};\]
		% 	\end{definition}
	
	We finish this subsection with an elementary statement about the pair $\fb_{\fsl_2}\subset \fsl_2$, which is the only non-formal input used below: it is applied in Lemma~\ref{lem:FiniteHomologicalDimension} and in the proof of Theorem~\ref{thm:IsomorphismInductionCoinduction}.
	
	\begin{lemma}\label{lem:Sl2InductionCoinduction}
		Let $\fsl_2=\langle e,h,f\rangle$, let $\fb_{\fsl_2}=\langle h,e \rangle$ be its Borel subalgebra and let $\Bbbk_{\mu \varepsilon}$ be the one dimensional $\fb_{\fsl_2}$-module of weight $\mu \varepsilon$, $\mu \in \mathbb{Z}$. Then
		\begin{enumerate}
			\item $\mathbf{L}^{(m)}\widetilde\Ind_{\fb_{\fsl_2}}^{\fsl_2}=0$ and $\mathbf{R}^{(m)}\widetilde\coInd_{\fb_{\fsl_2}}^{\fsl_2}=0$ for $m \geq 2$,
			\item the non-vanishing values on the one dimensional modules are
			\[\mathbf{L}^{(0)}\widetilde\Ind(\Bbbk_{\mu \varepsilon})=V_{\mu \varepsilon}~(\mu \geq 0),\qquad \mathbf{L}^{(1)}\widetilde\Ind(\Bbbk_{\mu \varepsilon})=V_{(-\mu-2) \varepsilon}~(\mu \leq -2),\]
			\[\mathbf{R}^{(0)}\widetilde\coInd(\Bbbk_{\mu \varepsilon})=V_{-\mu \varepsilon}~(\mu \leq 0),\qquad \mathbf{R}^{(1)}\widetilde\coInd(\Bbbk_{\mu \varepsilon})=V_{(\mu-2) \varepsilon}~(\mu \geq 2),\]
			\item consequently
			\begin{equation}\label{eq:Sl2InductionCoinduction}
				\Rcoind_{\fb_{\fsl_2}}^{\fsl_2}\simeq \Lind_{\fb_{\fsl_2}}^{\fsl_2}\circ\bigl(V(-2\varepsilon)\otimes \bullet\bigr)[-1],
			\end{equation}
			where $V(-2\varepsilon)$ is the one dimensional module of weight $-2\varepsilon$.
		\end{enumerate}
	\end{lemma}
	\begin{proof}
		Claims (1) and (2) are the Bott--Borel--Weil theorem for $\mathbb{P}^1$. On the level of the Grothendieck groups they are equivalent to
		\[[\widetilde\Ind_{\fb_{\fsl_2}}^{\fsl_2}]=(1-X^{-2})^{-1}(1-X^{-2}s),\qquad [\widetilde\coInd_{\fb_{\fsl_2}}^{\fsl_2}]=(1-X^{2})^{-1}(1-X^{2}s),\]
		the two expressions being obtained from one another by $X \mapsto X^{-1}$. Applying the first one to $X^{\mu}$ gives $\ch V_{\mu \varepsilon}$ for $\mu \geq 0$, zero for $\mu=-1$ and $-\ch V_{(-\mu-2)\varepsilon}$ for $\mu \leq -2$. Applying the second one gives $\ch V_{-\mu \varepsilon}$ for $\mu \leq 0$, zero for $\mu = 1$ and $-\ch V_{(\mu-2)\varepsilon}$ for $\mu \geq 2$.
		
		Comparing the two lines of (2) we get $\mathbf{R}^{(i)}\widetilde\coInd(\Bbbk_{\mu \varepsilon})\simeq \mathbf{L}^{(1-i)}\widetilde\Ind(\Bbbk_{(\mu-2) \varepsilon})$ for $i=0,1$. Since $\mathbf{L}^{(m)}$ sits in cohomological degree $-m$ and $\mathbf{R}^{(m)}$ in cohomological degree $m$, the passage from $1-i$ to $i$ is the shift by $[-1]$, and the twist $\mu \mapsto \mu-2$ is the tensoring by $V(-2\varepsilon)$. This is \eqref{eq:Sl2InductionCoinduction}.
	\end{proof}
	
	\subsection{Twist of a module by an automorphism}
	We need to introduce one more family of functors. Let $\fg$ be a Lie superalgebra and $\psi$ be an automorphism of $\fg$. Let $M$ be an $\fg$-module. The following definition is standard, see for example
	\begin{definition}\label{def:TwistedModule}
		$M^\psi$ is an $\fg$-module which is isomorphic to $M$ as a vector space and the action of $\fg$ is the following:
		\[X m^\psi:=(\psi(X)m)^\psi, ~\text{for }X \in \fg, m \in M.\]
	\end{definition}
	
	Then the map $\psi: M \mapsto M^{\psi}$ is an exact endofunctor of $\mathcal{O}(\fg)$.
	
	\subsection{Functors for the parabolic subalgebras of \texorpdfstring{$\Daff$}{Daff}}
	
	In this subsection we study the induction-restriction functors in the case of parabolic subalgebras of the Lie superalgebra $\Daff$. Consider the Lie superalgebras $\fp_{ij}\subset \fp_{123}$. Consider first the restriction functor
	$\Res_{\fp_{ij}}^{\fp_{123}}$. We first compute the action of this functor on the Grothendieck groups. Recall first that for $\{i,k\}=\{1,2,3\}\backslash\{j\}$
	\[K_0(\mathcal{O}^{gr}(\fp_{ik}))\simeq\mathbb{Z}[X_1^{\pm 1},X_2^{\pm 1},X_3^{\pm 1},t^{\pm 1}]^{s_j}((q^{\frac{1}{2}})),~K_0(\mathcal{O}^{gr}(\fp_{123}))\simeq\mathbb{Z}[X_1^{\pm 1},X_2^{\pm 1},X_3^{\pm 1},t^{\pm 1}]^{s_1,s_2,s_3}((q^{\frac{1}{2}})).\]  
	
	We use these isomorphisms to identify the Grothendieck groups with the corresponding subgroups of the rings of Laurent polynomials. So we get that the derived functors act on these rings of polynomials. In this subsection we explain these actions. 
	
	Note first that we have the natural inclusions
	\[[\Res_{\fp_{ik}}^{\fp_{123}}]:\mathbb{Z}[X_1^{\pm 1},X_2^{\pm 1},X_3^{\pm 1},t^{\pm 1}]^{s_1,s_2,s_3}((q^{\frac{1}{2}}))\hookrightarrow \mathbb{Z}[X_1^{\pm 1},X_2^{\pm 1},X_3^{\pm 1},t^{\pm 1}]^{s_j}((q^{\frac{1}{2}}))\]
	realizes the action of the restriction functor on the Grothendieck group. 
	
	Recall next that by Corollary~\ref{cor:ParabolicStable}
	the subalgebra $\fp_{ij}$ is stable under the automorphism $\pi_{(i,j)(k,0)}$, where $\{k\}=\{1,2,3\}\backslash \{i,j\}$. 
	\begin{lemma}
		The operator $[\pi_{(i,j)(k,0)}]$ acts on the ring of the Laurent polynomials by the following substitution:
		\[X_k \mapsto X_k, X_i \mapsto q^{\frac{1}{2}}X_i^{-1}, X_j \mapsto q^{\frac{1}{2}}X_j^{-1}.\]
	\end{lemma}
	\begin{proof}
		The twists over the automorphisms of a Lie superalgebra preserve the tensor products, i.e.\ for $M,N \in \mathcal{O}^{gr}(\fp_{ij})$ we have
		\[\pi_{(i,j)(k,0)}(M) \otimes \pi_{(i,j)(k,0)}(N)\simeq \pi_{(i,j)(k,0)}(M\otimes N).\]
		Therefore $[\pi_{(i,j)(k,0)}]$ is an automorphism of the polynomial ring. Moreover, $\pi_{(i,j)(k,0)}$ is an involutive automorphism and
		\[[\pi_{(i,j)(k,0)}](X^{\alpha_i})=X^{\alpha_j},~[\pi_{(i,j)(k,0)}](X^{\alpha_k})=X^{\alpha_0}.\]
		In other words,
		\begin{equation*}
			X_kX_iX_j^{-1} \longleftrightarrow  X_kX_i^{-1}X_j,~  X_k^{-1} X_iX_j\longleftrightarrow  qX_k^{-1} X_i^{-1}X_j^{-1}.
		\end{equation*}
		This allows us to compute the action of the operator on the variables.
	\end{proof}
	
	Now let us compute the action of the induction functor 
	$[\widetilde{\Ind}_{\fp_{ij}}^{\fp_{123}}]$ on the Grothendieck group.
	\begin{lemma}\label{lem:InductionOperator}
		\[[\widetilde{\Ind}_{\fp_{ij}}^{\fp_{123}}]=(1-X_{i}^{-2})^{-1}(1-X_{i}^{-2}s_i)(1-X_{j}^{-2})^{-1}(1-X_{j}^{-2}s_j)(1-tX^{-\alpha_k})(1-tX^{-\alpha_1-\alpha_2-\alpha_3}).\]
	\end{lemma}
	\begin{proof}
		Let us denote by $F_{ij}$ the operator 
		\begin{equation}\label{eq:FDefinition}
			F_{ij}:=(1-X_{i}^{-2})^{-1}(1-X_{i}^{-2}s_i)(1-X_{j}^{-2})^{-1}(1-X_{j}^{-2}s_j)(1-tX^{-\alpha_k})(1-tX^{-\alpha_1-\alpha_2-\alpha_3}).
		\end{equation}
		Then for $f \in \mathbb{Z}[X_1^{\pm 1},X_2^{\pm 1},X_3^{\pm 1},q^{\pm \frac{1}{2}},t^{\pm 1}]^{s_1,s_2,s_3}$ the operator $F_{ij}$ commutes with the multiplication by $f$:
		\[F_{ij} f =f F_{ij}.\]
		Moreover by Corollary~\ref{cor:CommutationMultiplicationSymmetric} one has
		the same property for the induction operator:
		\begin{equation*}
			[\widetilde{\Ind}_{\fp_{ij}}^{\fp_{123}}]f =f[\widetilde{\Ind}_{\fp_{ij}}^{\fp_{123}}].
		\end{equation*}
		
		Now let us compute both operators on the irreducible modules $V_{ij}(\mu)$ for $\mu_i,\mu_j\geq 2$. By skew-linearity it is enough to treat the modules without shifts. On the one hand by direct computation we get
		\begin{multline*}
			F_{ij}\left(X_i^{\mu_i}X_j^{\mu_j} \frac{X_k^{\mu_k}-X_k^{-\mu_k-2}}{1-X_k^{-2}}\right)\\=\ch(V_{123}(\mu))-t\ch(V_{123}(\mu-\alpha_k))-t\ch(V_{123}(\mu-\alpha_1-\alpha_2-\alpha_3))+t^2\ch(V_{123}(\mu-2\varepsilon_i-2\varepsilon_j)).
		\end{multline*}
		
		On the other hand let us compute the image of the induction functor $\LInd_{\fp_{ij}}^{\fp_{123}}$ on the irreducible module $V_{ij}(\mu)$.
		
		Consider the intermediate Lie superalgebra $\fp_{ij}^+:=\fp_{ij}\oplus \langle  e_{- \alpha_k}, e_{-\alpha_1-\alpha_2-\alpha_3}\rangle$, $\fp_{ij}\subset \fp_{ij}^+ \subset \fp_{123}$. Then by definition we have the following decomposition of the induction functor:
		
		\begin{equation}\label{eq:CompositionInductions}
			\widetilde{\Ind}_{\fp_{ij}}^{\fp_{123}}= \widetilde{\Ind}_{\fp_{ij}^+}^{\fp_{123}}\circ \Ind^{{\fp_{ij}^+}}_{\fp_{ij}}.
		\end{equation}
		
		The functor $ \Ind^{{\fp_{ij}^+}}_{\fp_{ij}}$ is exact. Therefore \eqref{eq:CompositionInductions} can be extended to the derived functors:
		\begin{equation*}\label{eq:DerivedCompositionInductions}
			\LInd_{\fp_{ij}}^{\fp_{123}}= \LInd_{\fp_{ij}^+}^{\fp_{123}}\circ \Ind^{{\fp_{ij}^+}}_{\fp_{ij}}.
		\end{equation*}
		
		Consider now an irreducible $\fp_{ij}$ module $V_{ij}(\mu)$. By definition of the induced module we have
		\begin{multline}\label{eq:FirstInductionCharacter}
			\ch  \Ind^{{\fp_{ij}^+}}_{\fp_{ij}}V_{ij}(\mu)= (1-tX^{-\alpha_k})(1-tX^{-\alpha_1-\alpha_2-\alpha_3})\ch V_{ij}(\mu)\\=(1-tX_i^{-1}X_j^{-1}X_k)(1-tX_1^{-1}X_2^{-1}X_3^{-1})\ch  V_{ij}(\mu).
		\end{multline}
		
		Assume now $\mu_i, \mu_j \geq 2$. The two generators added to $\fp_{ij}$ have weights $-\alpha_k$ and $-\alpha_1-\alpha_2-\alpha_3$, and $\alpha_k+\alpha_1+\alpha_2+\alpha_3=2\varepsilon_i+2\varepsilon_j$. Therefore the composition series of the module $\Ind^{{\fp_{ij}^+}}_{\fp_{ij}}V_{ij}(\mu)$ consists of four modules $V_{ij}(\mu)$, $V_{ij}(\mu-\alpha_k)\sh{0}{1}$, $V_{ij}(\mu- \alpha_1-\alpha_2-\alpha_3)\sh{0}{1}$, $V_{ij}(\mu-2\varepsilon_i-2\varepsilon_j)\sh{0}{2}$, in agreement with \eqref{eq:FirstInductionCharacter}. Since $\alpha_k$ and $\alpha_1+\alpha_2+\alpha_3$ both have coefficient $1$ at $\varepsilon_i$ and at $\varepsilon_j$, all these modules have dominant weights with respect to $h_i, h_j$. Therefore they are acyclic with respect to $\widetilde{\Ind}_{\fp_{ij}^+}^{\fp_{123}}$, therefore 
		\[\mathbf{L}^{>0}\widetilde{\Ind}_{\fp_{ij}^+}^{\fp_{123}}(\Ind^{{\fp_{ij}^+}}_{\fp_{ij}}V_{ij}(\mu))=0.\]
		However
		\[\widetilde{\Ind}_{\fp_{ij}^+}^{\fp_{123}}V_{ij}(\mu)=V_{123}(\mu)\]
		and by \eqref{eq:character3Irreducible} we have
		\[\ch \widetilde{\Ind}_{\fp_{ij}^+}^{\fp_{123}}V_{ij}(\mu)=(1-X_{i}^{-2})^{-1}(1-X_{i}^{-2}s_i)(1-X_{j}^{-2})^{-1}(1-X_{j}^{-2}s_j)\ch V_{ij}(\mu).\]
		
		Together with \eqref{eq:FirstInductionCharacter} this proves Lemma~\ref{lem:InductionOperator} for the irreducible representations $V_{ij}(\mu)$, $\mu_i,\mu_j \geq 2$. In other words,
		\[([\widetilde{\Ind}_{\fp_{ij}}^{\fp_{123}}]-F_{ij})\ch V_{ij}(\mu)=0.\]
		Moreover the operator $[\widetilde{\Ind}_{\fp_{ij}}^{\fp_{123}}]-F_{ij}$ commutes with the multiplication by the character of a representation of $\fp_{123}$. In particular, it commutes with $X_i+X_i^{-1}$ and $X_j+X_j^{-1}$. Note that 
		\[\ch V_{ij}(\mu+ \varepsilon_i)=\ch V_{ij}(\mu)X_i.\]
		Assume by induction that 
		\[([\widetilde{\Ind}_{\fp_{ij}}^{\fp_{123}}]-F_{ij})\ch V_{ij}(\mu)=0, \mu_i \geq \nu_i, \mu_j \geq \nu_j.\]
		Then 
		\[([\widetilde{\Ind}_{\fp_{ij}}^{\fp_{123}}]-F_{ij})(X_i+X_i^{-1})\ch V_{ij}(\mu)=0,\]
		and hence 
		\[([\widetilde{\Ind}_{\fp_{ij}}^{\fp_{123}}]-F_{ij})\ch V_{ij}(\mu)=0, \mu_i \geq \nu_i-1, \mu_j \geq \nu_j.\]
		Then by induction we get $[\widetilde{\Ind}_{\fp_{ij}}^{\fp_{123}}]=F_{ij}$.
	\end{proof}
	
	In order for the operator $[\widetilde{\Ind}_{\fp_{ij}}^{\fp_{123}}]$ used above to be defined at all, we need the induction functor to be of finite homological dimension, cf.\ Subsection~\ref{ssec:GrothendieckGroups}. This is the content of the following lemma.
	
	\begin{lemma}\label{lem:FiniteHomologicalDimension}
		$\mathbf{L}^{(m)}\widetilde{\Ind}_{\fp_{ij}}^{\fp_{123}}=0$ for $m \geq 3$. Consequently the operator
		\[[\widetilde{\Ind}_{\fp_{ij}}^{\fp_{123}}]=\sum_{m=0}^{2}(-1)^m\bigl[\mathbf{L}^{(m)}\widetilde{\Ind}_{\fp_{ij}}^{\fp_{123}}\bigr]\]
		is well defined on the Grothendieck groups.
	\end{lemma}
	\begin{proof}
		We use the factorization \eqref{eq:CompositionInductions} through the intermediate subalgebra $\fp_{ij}^+$. The functor $\Ind^{\fp_{ij}^+}_{\fp_{ij}}$ is exact, since $\U(\fp_{ij}^+)$ is free as a right $\U(\fp_{ij})$-module. Moreover $\fp_{ij}^+$ is obtained from $\fp_{ij}$ by adding two odd root vectors only, so no integrable quotient has to be taken. Hence it suffices to bound the homological dimension of $\widetilde{\Ind}_{\fp_{ij}^+}^{\fp_{123}}$.
		
		The complement of $\fp_{ij}^+$ in $\fp_{123}$ is spanned by the two even root vectors $e_{-2\varepsilon_i}=f_i$ and $e_{-2\varepsilon_j}=f_j$. Therefore $\widetilde{\Ind}_{\fp_{ij}^+}^{\fp_{123}}$ is the composition of the two integrable induction functors corresponding to the two commuting copies of $\fsl_2$ indexed by $i$ and $j$. By part (1) of Lemma~\ref{lem:Sl2InductionCoinduction} each of them has homological dimension one, so the composition has homological dimension at most two.
	\end{proof}
	\begin{remark}
		The number $2$ obtained here is exactly the shift appearing in Theorem~\ref{thm:IsomorphismInductionCoinduction}.
	\end{remark}
	
	In Section~\ref{sec:Orthogonality} the coinduction functor is applied to modules of the form $\Res_{\fp_{ij}}^{\fp_{123}}N$ with $N \in \mathcal{O}^{gr}(\fp_{123})$, that is, to objects of $\mathcal{O}^{gr}(\fp_{ij})$ and not of $\mathcal{O}^{\vee}(\fp_{ij})$. We therefore record separately that for our pair the coinduction is defined on $\mathcal{O}^{gr}$ as well.
	
	\begin{lemma}\label{lem:CoindPreservesO}
		$\widetilde\coInd_{\fp_{ij}}^{\fp_{123}}$ and its right derived functors map $\mathcal{O}^{gr}(\fp_{ij})$ to $\mathcal{O}^{gr}(\fp_{123})$, and $\mathbf{R}^{(m)}\widetilde\coInd_{\fp_{ij}}^{\fp_{123}}=0$ for $m \geq 3$.
	\end{lemma}
	\begin{proof}
		We again factor through $\fp_{ij}^+$. The complement of $\fp_{ij}$ in $\fp_{ij}^+$ is purely odd of dimension $(0|2)$, so that $\coInd^{\fp_{ij}^+}_{\fp_{ij}}M \simeq \Lambda(\fn)^*\otimes M$ is a finite direct sum of weight shifts of $M$. In particular this functor is exact and preserves both the boundedness of the weights and the integrability over the reductive part. For the second step, $\widetilde\coInd_{\fp_{ij}^+}^{\fp_{123}}$ is the Zuckerman type functor associated with the two commuting copies of $\fsl_2$ indexed by $i$ and $j$, and Lemma~\ref{lem:Sl2InductionCoinduction} gives both the required bound on the homological dimension and the fact that the values are integrable modules with weights bounded from above.
	\end{proof}
	
	\begin{definition}\label{def:Fdefinition}
		Let us introduce the functor 
		\[\mathcal{F}_{ij}:=\LInd_{\fp_{ij}}^{\fp_{123}} \circ \pi_{(i,j)(k,0)}\circ {\Res}_{\fp_{ij}}^{\fp_{123}}.\]
	\end{definition}
	
	The computations of this subsection immediately give us
	\begin{proposition}\label{prop:FEqual}
		\[[\mathcal F_{ij}]=F_{ij} \circ \pi_{(i,j)(k,0)}.\]
	\end{proposition}
    
	
	
	\section{Categorification of the genus two DAHA}\label{sec:Categorification}

    The genus two DAHA $\mathcal{A}_{q,t}$ was introduced in \cite{ArthamonovShakirov-2019} as the algebra generated by three multiplication operators $\hat{O}_{B_{ij}}=X_{ij}+X_{ij}^{-1}$ together with three commuting difference operators

    \begin{equation}\label{eq:OAOriginalForm}
		\hat O_{A_{ij}}=\hat O_{A_k}=\sum_{a,b\in\{\pm1\}}C^{a,b}_{i,j}\delta_{i}^a \delta_{j}^b,
	\end{equation}
	where
	\begin{equation}\label{eq:Cab}
		C^{a,b}_{i,j}=ab\,
		\frac{\bigl(1-t\,X_{k}X_{i}^{a}X_{j}^{b}\bigr)\bigl(1-t\,X_{k}^{-1}X_{i}^{a}X_{j}^{b}\bigr)}
		{t\,X_{i}^{a}X_{j}^{b}\bigl(X_{i}-X_{i}^{-1}\bigr)\bigl(X_{j}-X_{j}^{-1}\bigr)},
	\end{equation}
	and $\delta_{i}^a$ are the $q$-shift operators, i.e.\ the substitutions
	\[\delta_{i}^a(X_i)=q^{\frac{a}{2}} X_i.\]
    
    The operators $\hat{O}_{A_i}$ are the genus two analogue of the Macdonald operators, their basis of common eigenfunctions being the genus two Macdonald polynomials $\Psi_{j_1,j_2,j_3}$. The Mapping class group of $\Sigma_2$ acts by automorphisms on this algebra\cite{ArthamonovShakirov-2019}, while the specialization $t=q$ is isomorphic to the skein algebra $Sk_q(\Sigma_2)$ \cite{CookeSamuelson-2021}. In this section we construct a categorification of $\mathcal{A}_{q,t}$. Namely, the six generators lift to functors, the parameters $q,t$ arise from gradings and $\mathcal{A}_{q,t}$ is recovered on the Grothendieck group.
	
	Recall the  reflections from the Weyl group which  act by $s_{k}\colon X_{k}\mapsto X_{k}^{-1}$, fixing the two remaining $X$'s.
	
	\begin{lemma}
		Let $f$ be a symmetric function, $f =s_i(f)=s_j(f)$. Then 
		\[\delta_{i}^a \delta_{j}^bf=s_i^{\frac{a+1}{2}}s_j^{\frac{b+1}{2}}\pi_{(i,j)(k,0)}f.\]
	\end{lemma}
    \begin{proof}
	We have
	\[\pi s_i f = \pi f.\]
	Thus 
	\[\delta_{i}^a\delta_{j}^b f=s_i^{\frac{a+1}{2}}s_j^{\frac{b+1}{2}}\pi_{(i,j)(k,0)}s_i^{\frac{a-1}{2}}s_j^{\frac{b-1}{2}}f=s_i^{\frac{a+1}{2}}s_j^{\frac{b+1}{2}}\pi_{(i,j)(k,0)}f.\]
    %, whereas $\pi$ alone gives $q^{1/2}X_i^{-1}$, which equals $q^{-1/2}X_i$ after $s_i(f)=f$ is used.    
    \end{proof}
	
	Therefore, acting on symmetric functions, the operator \eqref{eq:OAOriginalForm} takes the form
	\begin{equation}\label{eq:OAOriginal}
		\hat O_{A_{ij}}=\hat O_{A_k}=\sum_{a,b\in\{\pm1\}}C^{a,b}_{i,j}\,s_{i}^{\frac{a+1}{2}}s_{j}^{\frac{b+1}{2}}\,\pi_{(i,j)(k,0)}.
	\end{equation}
	
	
	\begin{proposition}\label{prop:OAF}
		Recall the operator $F_{ij}$, \eqref{eq:FDefinition}. One has $\hat O_{A_{ij}}=t^{-1}F_{ij}\,\pi_{(i,j)(k,0)}$.
		% \begin{equation}\label{eq:F}
			% 	F_{ij}=\frac{1}{t}
			% 	\bigl(1-X_{k,i}^{-2}\bigr)^{-1}\bigl(1-X_{k,j}^{-2}\bigr)^{-1}
			% 	\bigl(1-X_{k,i}^{-2}s_{k,i}\bigr)\bigl(1-X_{k,j}^{-2}s_{k,j}\bigr)
			% 	\bigl(1-t\,X^{-\alpha_k}\bigr)
			% 	\bigl(1-t\,X^{-\alpha_i-\alpha_j-\alpha_k}\bigr).
			% \end{equation}
	\end{proposition}
	
	\begin{proof}
		Let us multiply both sides of the desired equality from the right by the Klein element $\pi_{(i,j)(k,0)}$. We get the following equality equivalent to the initial one:
		\begin{equation}\label{eq:core}
			\sum_{a,b\in\{\pm1\}}C^{a,b}_{i,j}\,
			s_{i}^{\frac{a+1}{2}}s_{j}^{\frac{b+1}{2}}=t^{-1}F_{ij}.
		\end{equation}
		
		Since $X_{i}-X_{i}^{-1}=X_{i}\bigl(1-X_{i}^{-2}\bigr)$, and similarly for $j$, the factor $\bigl(1-X_{i}^{-2}\bigr)^{-1} \bigl(1-X_{j}^{-2}\bigr)^{-1}$ is independent of $a,b$ and may be pulled out of the sum:
		\begin{equation}\label{eq:step1}
			\text{LHS of }\eqref{eq:core}=\frac{1}{t}\bigl(1-X_{i}^{-2}\bigr)^{-1}\bigl(1-X_{j}^{-2}\bigr)^{-1}\!\!\sum_{a,b\in\{\pm1\}}\!\! ab\,\frac{\bigl(1-t X_{k}X_{i}^{a}X_{j}^{b}\bigr)\bigl(1-t X_{k}^{-1}X_{i}^{a}X_{j}^{b}\bigr)} {X_{i}^{a+1}X_{j}^{b+1}}\,s_{i}^{\frac{a+1}{2}}s_{j}^{\frac{b+1}{2}}.
		\end{equation}
		Define now the group elements $\sigma_{a,b}:=s_{i}^{\frac{a+1}{2}}s_{j}^{\frac{b+1}{2}}$. If $a=1$ then $\sigma_{a,b}$ contains $s_{i}$ and sends $X_{i}^{a}=X_{i}$ to $X_{i}^{-1}$. If $a=-1$ then $\sigma_{a,b}$ fixes $X_{i}$ and the exponent is already $-1$. Either way $\bigl(X_{i}^{a}\bigr)^{\sigma_{a,b}}=X_{i}^{-1}$, and similarly for $X_{j}$. Hence for every $a,b\in\{\pm1\}$,
		\begin{align}
			\bigl(1-t X_{k}X_{i}^{a}X_{j}^{b}\bigr)
			\bigl(1-t X_{k}^{-1}X_{i}^{a}X_{j}^{b}\bigr)\sigma_{a,b}
			&=\sigma_{a,b}\bigl(1-t X_{k}X_{i}^{-1}X_{j}^{-1}\bigr)
			\bigl(1-t X_{k}^{-1}X_{i}^{-1}X_{j}^{-1}\bigr)
			\notag\\
			&=\sigma_{a,b}\bigl(1-t X^{-\alpha_k}\bigr)
			\bigl(1-t X^{-\alpha_1-\alpha_2-\alpha_3}\bigr),
			\label{eq:step2}
		\end{align}
		the last equality because
		$X_{k}X_{j}^{-1}X_{i}^{-1}=X^{-\alpha_k}$ and
		$X_{k}^{-1}X_{j}^{-1}X_{i}^{-1}=X^{-\alpha_i-\alpha_j-\alpha_k}$.
		The right-hand factors no longer depend on $a,b$, so they too leave the sum.
		
		Since $s_{i}$ fixes $X_{j}$ and $s_{j}$ fixes $X_{i}$, the sum
		factorizes, and for a single index
		\begin{equation}
			\sum_{a\in\{\pm1\}}a\,X_{i}^{-a-1}s_{i}^{\frac{a+1}{2}}
			=X_{i}^{-2}s_{i}-1=-\bigl(1-X_{i}^{-2}s_{i}\bigr).
		\end{equation}
		Multiplying the two factors, the signs cancel:
		\begin{equation}\label{eq:step3}
			\sum_{a,b\in\{\pm1\}}ab\,X_{i}^{-a-1}X_{j}^{-b-1}\,
			s_{i}^{\frac{a+1}{2}}s_{j}^{\frac{b+1}{2}}
			=\bigl(1-X_{i}^{-2}s_{i}\bigr)\bigl(1-X_{j}^{-2}s_{j}\bigr).
		\end{equation}
		Substituting \eqref{eq:step2} and \eqref{eq:step3} into \eqref{eq:step1}
		yields \eqref{eq:core}, which completes the proof.
	\end{proof}
	
	Then we get the following
	\begin{theoremB}\label{thm:MainTheorem}
		We have the following action of the functors $\mathcal{F}_{ij}$ on the Grothendieck group:
		\[[\mathcal{F}_{ij}]=t \hat O_{A_{ij}}.\]
		In other words, $\mathcal{F}_{ij}$ categorifies the genus two Macdonald operators.
	\end{theoremB}
  \begin{proof}
        By Proposition \ref{prop:FEqual} one has $[\mathcal{F}_{ij}]=F_{ij}\circ\pi_{(i,j)(k,0)}$, and by Proposition~\ref{prop:OAF} $\hat O_{A_{ij}}=t^{-1}F_{ij}\pi_{(i,j)(k,0)}$, whence $[\mathcal{F}_{ij}]=t\hat O_{A_{ij}}$. 
    \end{proof}
    
	\section{Orthogonality}\label{sec:Orthogonality}
	In this section we compute the $\Ext$ pairing of Subsection~\ref{ssec:ExtPairing} explicitly for the category $\mathcal{O}^{gr}(\fp_{123})$ and prove that the functors $\mathcal{F}_{ij}$ are self-adjoint with respect to it.
	
	\subsection{Projective modules}\label{ssec:ProjectiveModules}
	In this subsection we discuss the characters of projective modules to compute the pairing on the Grothendieck group in the explicit form.
	
	Let $R$ be the pro-nilpotent radical of the Lie superalgebra $\fp_{123}$, so that $\fp_{123}\simeq \fp_{123}^{red}\ltimes R$ with $\fp_{123}^{red}\simeq \fsl_2\oplus \fsl_2 \oplus \fsl_2$ spanned by $\{e_i,h_i,f_i\}$ in $z$-degree zero. Then
	\[R=\Bbbk K \oplus \spn\bigl\langle \{e_{\pm 2 \varepsilon_i}z^k,\ h_iz^k \mid k \geq 1\} \cup \{v_{\pm \varepsilon_1 }\otimes v_{\pm \varepsilon_2} \otimes v_{\pm \varepsilon_3} z^k \mid k \geq 0\}\bigr\rangle.\]
	Note that the central element $K$ has to be included in the radical: it is even and central, and it is not contained in $\fp_{123}^{red}$. We consider only modules of level zero, i.e.\ modules on which $K$ acts by zero. On such modules $\U(R)$ acts through $\U(R/\Bbbk K)$, which is why $K$ does not contribute to the character computed below. The three families $h_iz^k$, $k \geq 1$, are the root vectors of the imaginary roots $k\delta$, whose multiplicity is three, they produce the factor $(q;q)^{-3}$.

	Note that
	\[P_{123}(\mu) \simeq \U(\fp_{123})\otimes_{\U(\fp_{123}^{red})}V_{123}(\mu).\]
	
	Therefore there is an isomorphism of the graded vector spaces
	\[P_{123}(\mu) \simeq \U(R)\otimes V_{123}(\mu).\]
	
	Therefore we have the following equality of the characters:
	\[\ch (P_{123}(\mu))=\ch( \U(R))\ch ( V_{123}(\mu)).\]
	
	Note that
	\begin{equation*}
		\ch( \U(R))=\frac{\prod_{l \geq 0}\prod_{a,b,c\in\{\pm 1\}}(1-X_1^aX_2^bX_3^c q^lt)}{\prod_{l \geq 1}\prod_{i=1}^3(1-X_i^2 q^l)(1-X_i^{-2} q^l)}(q;q)^{-3},
	\end{equation*}
	where $(q;q):=\prod_{l \geq 1}(1-q^l)$ is the $q$-Pochhammer symbol.
	
	Thus we have
	\begin{equation*}
		\left(\frac{\prod_{l \geq 0}\prod_{a,b,c\in\{\pm 1\}}(1-X_1^aX_2^bX_3^c q^lt)}{\prod_{l \geq 1}\prod_{i=1}^3(1-X_i^2 q^l)(1-X_i^{-2} q^l)}[V_{123}(\mu)],[V_{123}(\mu')]\right)=\delta_{\mu,\mu'}(q;q)^{-3}.
	\end{equation*}
	
	For $f \in \mathbb{Z}[X_1^{\pm 1},X_2^{\pm 1},X_3^{\pm 1},t^{\pm 1}]((q^{\pm \frac{1}{2}}))$ we denote by $f_0 \in \mathbb{Z}[t^{\pm 1}]((q^{\pm \frac{1}{2}}))$ its $X$-free term. Then for the modules over a finite-dimensional semisimple Lie algebra $\fg$ we have
	\[\left(\prod_{\alpha \in \Phi_{\fg}}(1-X^\alpha)\ch(V_\lambda)\ch(V_{\lambda'})\right)_0=\delta_{\lambda,\lambda'}.\]
	Here $\Phi_\fg$ is the root system of the Lie algebra $\fg$. Note that $\fp_{123}^{red}\simeq \fsl_2 \oplus \fsl_2 \oplus \fsl_2$.
	
	Thus
	\[\left(\prod_{i=1}^3(1-X_i^2)(1-X_i^{-2})\ch V_{123}(\mu)\ch V_{123}(\mu')\right)_0=\delta_{\mu,\mu'}.\]
	
	This equality can be rephrased in the following way:
	\begin{equation}\label{eq:BilinearFormOnBasis}
		(q;q)^3\left(\frac{\prod_{l \geq 0}\prod_{i=1}^3(1-X_i^2 q^l)(1-X_i^{-2} q^l)}{\prod_{l \geq 0}\prod_{a,b,c\in\{\pm 1\}}(1-X_1^aX_2^bX_3^c q^lt)}[P_{123}(\mu)][V_{123}(\mu')]\right)_0=\delta_{\mu,\mu'}.
	\end{equation}
	
	Let us define
	\[\Delta:=(q;q)^3\frac{\prod_{l \geq 0}\prod_{i=1}^3(1-X_i^2 q^l)(1-X_i^{-2} q^l)}{\prod_{l \geq 0}\prod_{a,b,c\in\{\pm 1\}}(1-X_1^aX_2^bX_3^c q^lt)}.\]
	
	For $f \in \mathbb{Z}[X_1^{\pm 1},X_2^{\pm 1},X_3^{\pm 1},t^{\pm 1}]((q^{\frac{1}{2}}))$ we denote by $f^\star$ the substitution
	\[f^\star:=f|_{t \mapsto t^{-1},\,q \mapsto q^{-1}}.\]
	
	Then, extending the equality \eqref{eq:BilinearFormOnBasis} by skew-linearity, one gets the following lemma.
	
	\begin{lemma}\label{lem:ExtPairingExplicit}
		The $\Ext$-pairing on the Grothendieck ring can be expressed in the following form:
		\[(f,g)=(f \Delta g^\star)_0.\]
	\end{lemma}
	
	\subsection{Orthogonality of operators}
	Throughout this subsection $\{i,j,k\}=\{1,2,3\}$, we abbreviate $\pi:=\pi_{(i,j)(k,0)}$, and $M,N$ are objects of $\mathcal{O}^{gr}(\fp_{123})$. Recall that by Lemma~\ref{lem:CoindPreservesO} all the functors below are defined on $\mathcal{O}^{gr}$.
	
	By Proposition~\ref{prop:Adjoint} the functor ${\Lind}_{\fp_{ij}}^{\fp_{123}}$ is left adjoint to $\Res_{\fp_{ij}}^{\fp_{123}}$, so that for every $n \in \mathbb{Z}$
	\[\Hom_{gr}(M[n],\Res_{\fp_{ij}}^{\fp_{123}}N)\simeq
	\Hom_{gr}({\Lind}_{\fp_{ij}}^{\fp_{123}}M[n],N).\]
	Recall also that $\pi$ is an involution, hence for the twist functor of Definition~\ref{def:TwistedModule} one has $\Hom(A^{\pi},B)\simeq \Hom(A,B^{\pi})$: an $\fp_{ij}$-morphism $A^\pi \rightarrow B$ is the same thing as an $\fp_{ij}$-morphism $A \rightarrow B^{\pi^{-1}}=B^{\pi}$. Combining the adjunction for $\Lind$, this twist and the adjunction for $\Rcoind$ we get
	\begin{multline*}
		\Hom_{gr}({\Lind}_{\fp_{ij}}^{\fp_{123}}\circ \pi \circ \Res_{\fp_{ij}}^{\fp_{123}}M[n],N)\simeq  \Hom_{gr}(\pi \circ \Res_{\fp_{ij}}^{\fp_{123}} M[n],\Res_{\fp_{ij}}^{\fp_{123}} N)\\ \simeq 
		\Hom_{gr}(\Res_{\fp_{ij}}^{\fp_{123}}M[n],\pi \circ \Res_{\fp_{ij}}^{\fp_{123}} N)\simeq  \Hom_{gr}(M[n],{\Rcoind}_{\fp_{ij}}^{\fp_{123}}\circ \pi \circ \Res_{\fp_{ij}}^{\fp_{123}} N).
	\end{multline*}
	
	Therefore the following operators on the Grothendieck group are adjoint:
	\begin{equation}\label{eq:AdjointOperators}
		(t \hat O_{A_{ij}}[M],[N])=([{\Lind}_{\fp_{ij}}^{\fp_{123}}\circ \pi \circ \Res_{\fp_{ij}}^{\fp_{123}}][M],[N])= 
		([M],[{\Rcoind}_{\fp_{ij}}^{\fp_{123}}\circ \pi \circ \Res_{\fp_{ij}}^{\fp_{123}}][N]).
	\end{equation}
	
	It remains to identify the operator on the right hand side. It is possible to compute it directly, but we prove a somewhat stronger claim on the level of functors.
	\begin{theoremA}\label{thm:IsomorphismInductionCoinduction}
		There is the following isomorphism of derived functors:
		\[{\Rcoind}_{\fp_{ij}}^{\fp_{123}}\circ \pi \circ \Res_{\fp_{ij}}^{\fp_{123}}\simeq {\Lind}_{\fp_{ij}}^{\fp_{123}}\circ \pi \circ 	\Res_{\fp_{ij}}^{\fp_{123}}\sh{0}{-2}[-2].\]
	\end{theoremA}
	
	\begin{proof}
		It is enough to prove that 
		\begin{equation}\label{eq:IndCoindComparison}
			{\Rcoind}_{\fp_{ij}}^{\fp_{123}}\simeq {\Lind}_{\fp_{ij}}^{\fp_{123}}\sh{0}{-2}[-2].
		\end{equation}
		
		{\it Step 1: factorization.} Recall the intermediate superalgebra $\fp_{ij}^+$ of the proof of Lemma~\ref{lem:InductionOperator}. The functors ${\Ind}_{\fp_{ij}}^{\fp_{ij}^+}$ and ${\coInd}_{\fp_{ij}}^{\fp_{ij}^+}$ are exact, because the quotient $\fn:=\fp_{ij}^+/\fp_{ij}$ is purely odd of dimension $(0|2)$ and no integrable quotient or submodule has to be taken. Therefore
		\begin{equation}\label{eq:FactorizationOfBoth}
			{\Lind}_{\fp_{ij}}^{\fp_{123}}\simeq {\Lind}_{\fp_{ij}^+}^{\fp_{123}}\circ {\Ind}_{\fp_{ij}}^{\fp_{ij}^+},\qquad {\Rcoind}_{\fp_{ij}}^{\fp_{123}} \simeq {\Rcoind}_{\fp_{ij}^+}^{\fp_{123}}\circ {\coInd}_{\fp_{ij}}^{\fp_{ij}^+}.
		\end{equation}
		
		{\it Step 2: the odd part.} Note that $\fn$ is a two dimensional $\fp_{ij}$-module (it is the quotient $\fp_{ij}^+/\fp_{ij}$ which is a module. The span of $e_{-\alpha_k}$ and $e_{-\alpha_1-\alpha_2-\alpha_3}$ inside $\fp_{ij}^+$ is not a $\fp_{ij}$-submodule, for instance $[e_{2\varepsilon_i},e_{-\alpha_k}]$ is proportional to $e_{\alpha_j}$). Its top exterior power is one dimensional of weight $-\alpha_k-\alpha_1-\alpha_2-\alpha_3=-2\varepsilon_i-2\varepsilon_j$ and of odd degree $2$, i.e.\ $\Lambda^{top}\fn \simeq V_{ij}(-2\varepsilon_i-2\varepsilon_j)\sh{0}{2}$. By the Poincar\'e--Birkhoff--Witt theorem $\U(\fp_{ij}^+)\simeq \U(\fp_{ij})\otimes \Lambda(\fn)$, whence
		\begin{equation} \label{eq:OddInduction}
			{\Ind}_{\fp_{ij}}^{\fp_{ij}^+}(\bullet)\simeq {\coInd}_{\fp_{ij}}^{\fp_{ij}^+}\bigl(\bullet \otimes V_{ij}(-2\varepsilon_i-2\varepsilon_j)\sh{0}{2}\bigr).
		\end{equation}
		Indeed, \eqref{eq:OddInduction} follows from the isomorphisms of $\fh$-modules ${\Ind}_{\fp_{ij}}^{\fp_{ij}^+}M \simeq \Lambda(\fn)\otimes M$ and ${\coInd}_{\fp_{ij}}^{\fp_{ij}^+}M \simeq \Lambda(\fn)^*\otimes M$, together with $\Lambda(\fn)\simeq \Lambda(\fn)^*\otimes \Lambda^{top}\fn$. Equivalently,
		\begin{equation}\label{eq:OddCoinduction}
			{\coInd}_{\fp_{ij}}^{\fp_{ij}^+}(\bullet)\simeq {\Ind}_{\fp_{ij}}^{\fp_{ij}^+}\bigl(\bullet \otimes V_{ij}(2\varepsilon_i+2\varepsilon_j)\sh{0}{-2}\bigr).
		\end{equation}
		
		{\it Step 3: the even part.} The complement of $\fp_{ij}^+$ in $\fp_{123}$ is spanned by $f_i$ and $f_j$, so that $\widetilde\Ind_{\fp_{ij}^+}^{\fp_{123}}$ and $\widetilde\coInd_{\fp_{ij}^+}^{\fp_{123}}$ are the functors attached to the two commuting copies of $\fsl_2$ indexed by $i$ and $j$. Applying \eqref{eq:Sl2InductionCoinduction} twice we get
		\begin{equation}\label{eq:EvenIndCoind}
			{\Rcoind}_{\fp_{ij}^+}^{\fp_{123}} \simeq {\Lind}_{\fp_{ij}^+}^{\fp_{123}}\circ \bigl(V(-2\varepsilon_i-2\varepsilon_j)\otimes \bullet \bigr)[-2].
		\end{equation}
		
		{\it Step 4: the two twists cancel.} The modules $V_{ij}(\pm(2\varepsilon_i+2\varepsilon_j))$ are one dimensional, so tensoring by them commutes with induction. Substituting \eqref{eq:OddCoinduction} and \eqref{eq:EvenIndCoind} into the second identity of \eqref{eq:FactorizationOfBoth} we obtain
		\begin{multline*}
			{\Rcoind}_{\fp_{ij}}^{\fp_{123}}(\bullet)\simeq {\Lind}_{\fp_{ij}^+}^{\fp_{123}}\Bigl({\Ind}_{\fp_{ij}}^{\fp_{ij}^+}\bigl(\bullet \otimes V(2\varepsilon_i+2\varepsilon_j)\sh{0}{-2}\bigr)\otimes V(-2\varepsilon_i-2\varepsilon_j)\Bigr)[-2]\\
			\simeq {\Lind}_{\fp_{ij}^+}^{\fp_{123}}{\Ind}_{\fp_{ij}}^{\fp_{ij}^+}\bigl(\bullet \otimes V(2\varepsilon_i+2\varepsilon_j)\otimes V(-2\varepsilon_i-2\varepsilon_j)\bigr)\sh{0}{-2}[-2]\simeq {\Lind}_{\fp_{ij}}^{\fp_{123}}(\bullet)\sh{0}{-2}[-2],
		\end{multline*}
		which is \eqref{eq:IndCoindComparison}.
	\end{proof}
	
	\begin{corollary}\label{cor:CoindOperator}
		\[[{\Rcoind}_{\fp_{ij}}^{\fp_{123}}\circ \pi \circ \Res_{\fp_{ij}}^{\fp_{123}}]=t^{-1} \hat O_{A_{ij}}.\]
	\end{corollary}
	\begin{proof}
		By Theorem~\ref{thm:IsomorphismInductionCoinduction} the left hand side equals $(-t)^{-2}(-1)^{-2}[\mathcal{F}_{ij}]$, the first factor coming from the shift $\sh{0}{-2}$ by \eqref{eq:ShiftDefinition} and the second one from the homological shift $[-2]$. Both equal $t^{-2}$ and $1$ respectively. Since $[\mathcal{F}_{ij}]=t\hat O_{A_{ij}}$ by Theorem~\ref{thm:MainTheorem}, the claim follows.
	\end{proof}
	
	\begin{theoremC}\label{thm:SelfAdjointness}
		The operators $\hat O_{A_{ij}}$ are self-adjoint with respect to the pairing $(-,-)$ on $K_0(\mathcal{O}^{gr}(\fp_{123}))$:
		\begin{equation*}
			(\hat O_{A_{ij}}f,g)=(f,\hat O_{A_{ij}}g).
		\end{equation*}
	\end{theoremC}
	\begin{proof}
		Combining \eqref{eq:AdjointOperators} with Corollary~\ref{cor:CoindOperator} we get $(t\hat O_{A_{ij}}f,g)=(f,t^{-1}\hat O_{A_{ij}}g)$. The pairing is $\mathbb{Z}[t^{\pm 1}]((q^{\pm \frac{1}{2}}))$-linear in the first argument and skew-linear in the second one, so the left hand side equals $t(\hat O_{A_{ij}}f,g)$ and the right hand side equals $t(f,\hat O_{A_{ij}}g)$. Cancelling $t$ gives the claim.
	\end{proof}
	
	\begin{corollary}\label{cor:Orthogonality}
		Let $\Psi_{i_1,i_2,i_3}$ and $\Psi_{j_1,j_2,j_3}$ be rank $2$ Macdonsld polynomials. If $(i_1,i_2,i_3) \neq (j_1,j_2,j_3)$, then they are orthogonal, i. e. 
        \[(\Psi_{i_1,i_2,i_3},\Psi_{j_1,j_2,j_3})=0.\]
	\end{corollary}
	\begin{proof}
		By Theorem~\ref{thm:SelfAdjointness}, $\lambda_{ij}(f,g)=(\hat O_{A_{ij}}f,g)=(f,\hat O_{A_{ij}}g)=\mu_{ij}^{\star}(f,g)$, the last equality by the skew-linearity of the pairing in the second argument.
	\end{proof}
	
    \section*{Acknowledgements}
    S.A. would like to thank Shamil Shakirov for a multi-year collaboration that resulted in \cite{ArthamonovShakirov-2019} that became a prerequisite for the current manuscript. The potential connection between \cite{ArthamonovShakirov-2019} and the combinatorics of the root system of $D(2|1,\alpha)$ was brought to the attention of S.A. by Junichi Shiraishi during his visit to BIMSA in December 2024. The work of S.A. was supported by a grant of the Beijing Natural Science Foundation IS-25025.
    
    \textit{AI statement:} A large variety of AI tools was used during all stages of work on the project. DeepSeek v3.2 and v4, Gemini 3.1 Pro, Kimi K2.6 and K3, and Claude Opus 4.8, Opus 5 and Fable 5 were used as coding assistants in C++ and Mathematica for calculations that resulted in the starting character formulas. Claude Opus 5, Claude Fable 5 and Kimi K3 were used while working on the manuscript, but only to improve the clarity and flow of text that had already been written manually. In particular, the TikZ code for Figure \ref{fig:RootSystem} was produced by K3 based on the blackboard photograph and refined by Claude Opus 5. All proofs and logical content were inserted manually. Before finalizing the manuscript, the use of AI tools was suspended, so that the final version of the text was proofread and edited by the authors alone.
    
	\printbibliography
	
\end{document}